\documentclass{article}


% if you need to pass options to natbib, use, e.g.:
\PassOptionsToPackage{numbers, compress}{natbib}
% before loading neurips_2025


% ready for submission
% \usepackage{neurips_2025}


% to compile a preprint version, e.g., for submission to arXiv, add add the
% [preprint] option:
%     \usepackage[preprint]{neurips_2025}


% to compile a camera-ready version, add the [final] option, e.g.:
\usepackage[final]{neurips_2025}


% to avoid loading the natbib package, add option nonatbib:
%    \usepackage[nonatbib]{neurips_2025}


\usepackage[utf8]{inputenc} % allow utf-8 input
\usepackage[T1]{fontenc}    % use 8-bit T1 fonts
\usepackage{hyperref}       % hyperlinks
\usepackage{url}            % simple URL typesetting
\usepackage{booktabs}       % professional-quality tables
\usepackage{amsfonts}       % blackboard math symbols
\usepackage{nicefrac}       % compact symbols for 1/2, etc.
\usepackage{microtype}      % microtypography
\usepackage{xcolor}         % colors


% 新加的包
\usepackage{graphicx}
%\usepackage{subfigure}
\usepackage{subcaption}
\usepackage{algorithm}
\usepackage{algorithmic}
\usepackage{amsmath}
\usepackage{amssymb}
\usepackage{mathtools}
\usepackage{amsthm}
\usepackage{framed}
\usepackage{multirow}
\usepackage{svg}
\usepackage{enumitem}
\usepackage{wrapfig}
\usepackage{booktabs}

\usepackage[export]{adjustbox}

\usepackage[capitalize,noabbrev]{cleveref}
\usepackage{CJKutf8}


\newcommand{\ns}[1]{{\color{orange}{[[\textbf{ns: }#1]]}}}
\newcommand{\cx}[1]{{\color{green}{[[\textbf{cx: }#1]]}}}


\makeatletter
\newcommand{\samethanks}[1][\value{footnote}]{\footnotemark[#1]}
\makeatother


\input{math_commands}


\title{Large Language Diffusion Models}


% The \author macro works with any number of authors. There are two commands
% used to separate the names and addresses of multiple authors: \And and \AND.
%
% Using \And between authors leaves it to LaTeX to determine where to break the
% lines. Using \AND forces a line break at that point. So, if LaTeX puts 3 of 4
% authors names on the first line, and the last on the second line, try using
% \AND instead of \And before the third author name.



\author{%
  Shen Nie$^{1,2,3}$\thanks{Equal contribution.}~~\thanks{Work done during an internship at Ant Group.} \quad
  Fengqi Zhu$^{1,2,3}$\samethanks[1]~~\samethanks[2] \quad
  Zebin You$^{1,2,3}$\samethanks[2] \quad
  Xiaolu Zhang$^{4}$\thanks{Project leaders.} \quad
  Jingyang Ou$^{1,2,3}$ \\
  \textbf{Jun Hu}$^{4}$\samethanks[3] \quad
  \textbf{Jun Zhou}$^{4}$ \quad
  \textbf{Yankai Lin}$^{1,2,3}$\samethanks[3] \quad
  \textbf{Ji-Rong Wen}$^{1,2,3}$ \quad
  \textbf{Chongxuan Li}$^{1,2,3}$\samethanks[3]~~\thanks{Correspondence to Chongxuan Li.} \\
  $^1$ Gaoling School of Artificial Intelligence, Renmin University of China \\
  $^2$ Beijing Key Laboratory of Research on Large Models and Intelligent Governance \\
  $^3$ Engineering Research Center of Next-Generation Intelligent Search and Recommendation, MOE \\
  $^4$ Ant Group\\
  \texttt{\{nieshen,fengqizhu,chongxuanli\}@ruc.edu.cn} 
}





\begin{document}


\maketitle

\begin{abstract}
  The capabilities of large language models (LLMs) are widely regarded as relying on autoregressive models (ARMs). We challenge this notion by introducing \emph{LLaDA}, a diffusion model trained from scratch under the pre-training and supervised fine-tuning (SFT) paradigm. LLaDA employs a forward data masking process and a reverse generation process, parameterized by a Transformer to predict masked tokens. It provides a principled generative approach for probabilistic inference by optimizing a likelihood lower bound. Across extensive benchmarks on general tasks, math, code, and so on, LLaDA demonstrates strong \emph{scalability} and performs comparably to our self-constructed ARM baselines. Remarkably, LLaDA 8B is competitive with strong LLMs like LLaMA3 8B in \emph{in-context learning} and, after SFT, exhibits impressive \emph{instruction-following} abilities in case studies such as multi-turn dialogue. Moreover, LLaDA addresses the reversal curse, surpassing GPT-4o in a reversal poem completion task. Our findings show the promise of diffusion models for language modeling at scale and challenge the common assumption that core LLM capabilities discussed above inherently depend on ARMs. Project page and codes: \url{https://ml-gsai.github.io/LLaDA-demo/}.
\end{abstract}

\section{Introduction}
\label{sec:introduction}


Large language models (LLMs)~\citep{zhao2023survey} fall entirely within the framework of generative modeling. Specifically, LLMs aim to capture the true but unknown language distribution \( p_{\textrm{data}}(\cdot) \) by optimizing a model distribution \( p_{\theta}(\cdot) \) through maximum likelihood estimation, or equivalently KL divergence minimization between the two distributions:
\begin{align}
\label{eq:llm}
 \underbrace{\max_{\theta} \mathbb{E}_{ p_{\textrm{data}}(x) }\log p_{\theta}(x) \Leftrightarrow \min_{\theta} \textrm{KL}(p_{\textrm{data}}(x) || p_{\theta}(x))}_{\textrm{Generative modeling principles}}.
\end{align}


The predominant approach relies on the autoregressive modeling (ARM)—commonly referred to as the ``next-token prediction'' paradigm—to define the model distribution:
\begin{align}
\label{eq:autoregressive}
\underbrace{p_{\theta}(x) = p_\theta(x^1)   \prod_{i=2}^L p_\theta(x^i \mid x^{1}, \dots, x^{i-1})}_{\textrm{Autoregressive formulation}},
\end{align}
where \( x \) is a sequence of length \( L \), and \( x^i \) is the \( i \)-th token. This paradigm has proven remarkably effective~\citep{radford2018improving,radford2019language,brown2020language,chatgpt} and has become the foundation of current LLMs. Despite its widespread adoption, a fundamental question remains unanswered: \textit{Is the autoregressive paradigm the only path to achieving the core capabilities of LLMs, such as scalability, in-context learning, and instruction-following?}  


\begin{figure}[t!]
  \centering
  \begin{subfigure}{0.46\textwidth}
    \centering
    \includegraphics[width=\linewidth]{imgs/LLaDA_vs_LLaMA.pdf}
  \end{subfigure}\hfill
  \begin{subfigure}{0.49\textwidth}
    \centering
    \includegraphics[width=\linewidth]{imgs/LLaDA_vs_LLaMA_chat.pdf}
  \end{subfigure}
  \caption{\textbf{Zero/Few‑Shot Benchmarks.} We scale LLaDA to 8B parameters from scratch and observe competitive zero/few‑shot performance compared with strong autoregressive LLMs~\citep{dubey2024llama}.}
  \label{fig:LLaDA_vs_LLaMA}
\end{figure}




We argue that the answer is \emph{not} a simple ``yes''. The key insight overlooked previously is:
It is the \emph{generative modeling principles} (i.e., Eq. (\ref{eq:llm})), \emph{rather than the autoregressive formulation} (i.e., Eq. (\ref{eq:autoregressive})) itself, that fundamentally underpin the essential properties of LLMs.

In particular, we argue that \emph{scalability} is primarily a consequence of the interplay between Transformers~\citep{vaswani2017attention}, model size, data size, and \emph{Fisher consistency}\footnote{It suggests the ability to recover the true data distribution with infinite data, a sufficiently large network and optimal training.}~\citep{fisher1922mathematical} induced by the generative principles in Eq.~(\ref{eq:llm}), rather than a unique result of the ARMs in Eq.~(\ref{eq:autoregressive}). The success of diffusion transformers~\citep{bao2023all,peebles2023scalable} on visual data~\citep{videoworldsimulators2024} supports this claim. Furthermore, the \emph{instruction-following} and \emph{in-context learning}~\citep{brown2020language} capabilities appear to be intrinsic properties of all conditional generative models on structurally consistent linguistic tasks, rather than exclusive advantages of ARMs. In addition, while ARMs can be interpreted as a \emph{lossless data compressor}~\citep{deletanglanguage,huang2024compression}, any sufficiently expressive probabilistic model can achieve similar capabilities~\citep{shannon1948mathematical}.

However, certain inherent limitations of LLMs can be directly attributed to their autoregressive nature. For instance, the left-to-right generation process restricts their ability to handle reversal reasoning tasks~\citep{berglund2023reversal}, highlighting a representative failure in the generalization capabilities of current models.
 
Motivated by these insights, we introduce \emph{LLaDA (Large Language Diffusion with mAsking)} to investigate whether the capabilities exhibited by LLMs can emerge from generative modeling principles beyond ARMs, thereby addressing the fundamental question posed earlier. In contrast to traditional ARMs, LLaDA leverages a masked diffusion model (MDM)~\citep{austin2021structured,lou2023discrete,shi2024simplified,sahoo2024simple,ou2024your}, which incorporates a forward data masking process and trains a \emph{mask predictor} to approximate its reverse process. This design enables LLaDA to construct a model distribution with bidirectional dependencies and optimize a variational lower bound of its log-likelihood, offering a principled and previously unexplored perspective on the core capabilities of LLMs discussed above.


We adopt the standard pipeline of data preparation, pre-training, supervised fine-tuning (SFT), and evaluation, scaling LLaDA to an unprecedented language diffusion of size 8B. In particular, LLaDA 8B was pre-trained from scratch on 2.3 trillion tokens using 0.13 million H800 GPU hours, followed by SFT on 4.5 million pairs. Across diverse tasks, including language understanding, math, code, and Chinese, LLaDA demonstrates the following contributions:
\begin{itemize}[leftmargin=*, itemsep=0pt]
    \item LLaDA scales effectively to a compute budget of $10^{23}$ FLOPs, achieving comparable results to ARM baselines trained on the same data across six tasks, e.g., MMLU and GSM8K.
    \item The pre-trained LLaDA 8B Base surpasses LLaMA2 7B Base~\citep{touvron2023llama2} on nearly all 15 standard zero/few-shot learning tasks while performing on par with LLaMA3 8B Base~\citep{dubey2024llama}, showcasing effective in-context learning capability.
    \item LLaDA significantly enhances the ability to follow instructions after SFT, as demonstrated in case studies such as multi-turn dialogue.
    \item LLaDA effectively breaks the reversal curse~\citep{berglund2023reversal} with consistent performance across forward and reversal tasks. Notably, it outperforms GPT-4o in a reversal poem completion task. 

\end{itemize}


% \cx{citation. polish change. summary of contributions}

% Overall, LLaDA scales up to a compute budget of $10^{23}$ FLOPs and achieves performance comparable to our self-constructed ARM baselines trained on the same data across six diverse tasks, such as MMLU and GSM8K, demonstrating strong \emph{scalability} competitive with ARMs. Further, the pre-trained LLaDA 8B Base outperforms LLaMA2 7B Base~\citep{touvron2023llama2} on most of 15 standard zero/few-shot benchmarks, and performs competitively with LLaMA3 8B Base~\citep{dubey2024llama}, showcasing effective \emph{in-context learning}. After SFT, LLaDA 8B Instruct exhibits robust \emph{instruction-following} ability, as evidenced by case studies such as multi-turn dialogue interactions. Finally, LLaDA demonstrates strong \emph{reversal reasoning} capabilities, achieving consistent performance on both forward and reversal tasks, and even surpassing GPT-4o in a reversal poem completion benchmark.



\begin{figure*}[t!]
    \centering
    \includegraphics[width=0.98\linewidth]{imgs/overview.pdf}
    \vspace{-.15cm}
    \caption{\textbf{Overview of LLaDA.} (a) Pre-training. LLaDA is trained on text with random masks applied independently to all tokens at the same ratio $t \sim U[0, 1]$. (b) SFT. Only response tokens are possibly masked. (c) Sampling. LLaDA simulates a diffusion process from $t = 1$ (fully masked) to $t = 0$ (unmasked), predicting all masks simultaneously at each step with flexible remask strategies.}
    \label{fig:overview}
    \vspace{-.15cm}
\end{figure*}



\section{Approach}
\label{sec:approach}

In this section, we introduce the probabilistic formulation\footnote{Here, we focus on the approach of LLaDA. A rigorous formulation of MDM is provided in Appendix~\ref{app:formulation} for interested readers.}, along with the pre-training, supervised fine-tuning, and inference procedures for LLaDA, as illustrated in Fig.~\ref{fig:overview}.


\subsection{Probabilistic Formulation}

Unlike ARMs in Eq.~(\ref{eq:autoregressive}), LLaDA defines a model distribution \( p_\theta(x_0) \) through a \emph{forward process} and a \emph{reverse process}~\citep{austin2021structured,lou2023discrete,shi2024simplified,sahoo2024simple,ou2024your}. The forward process gradually masks tokens independently in \( x_0 \) until the sequence is fully masked at \( t = 1 \). For \( t \in (0, 1) \), the sequence \( x_t \) is partially masked, with each being masked with probability \( t \) or remaining unmasked with probability \( 1 - t \). The reverse process recovers the data distribution by iteratively predicting masked tokens as \( t \) moves from \( 1 \) to \( 0 \).

The core of LLaDA is a \emph{mask predictor}, a parametric model \( p_\theta(\cdot|x_t) \) that takes \( x_t \) as input and predicts all masked tokens (denoted as \textrm{M}) simultaneously. It is trained using a cross-entropy loss computed only on the masked tokens~\citep{shi2024simplified,sahoo2024simple,ou2024your}:
\begin{align}
\label{eq:objective}
   \mathcal{L}(\theta)  \triangleq   -  \mathbb{E}_{t, x_0,  x_t} \left[\frac{1}{t} \sum_{ i = 1 }^L \textbf{1}[x_t^i = \textrm{M}] \log p_{\theta}(x_0^i|x_t) \right] , 
\end{align}
where \( x_0 \) is a training sample, \( t \) is a continuous random variable drawn uniformly from \( [0, 1] \), \( x_t \) is sampled from the forward process and $L$ is the sequence length. The indicator function \( \textbf{1}[\cdot] \) ensures that the loss is computed only for masked tokens.

Once trained, we can simulate a reverse process (see Sec.~\ref{sec:inference} for details) parameterized by the mask predictor and define the model distribution \( p_\theta(x_0) \) as the marginal distribution induced at \( t = 0 \). The loss function in Eq.~(\ref{eq:objective}) has been proven to be an upper bound on the negative log-likelihood of the model distribution, making it a principled objective for generative modeling:
\begin{align}
\label{eq:bound}
    - \mathbb{E}_{p_{\textrm{data}}(x_0)} \left[\log p_\theta(x_0) \right]  \le  \mathcal{L}(\theta).
\end{align} 

Notably, LLaDA employs a masking ratio that varies randomly between 0 and 1 while BERT~\citep{devlin2018bert} uses a fixed ratio. The subtle differences have significant implications, especially at scale: as shown in Eq.~(\ref{eq:bound}), LLaDA is a principled generative model with the potential to perform in-context learning and instruction-following naturally, akin to LLMs. Moreover, its generative perspective implies strong scalability with large data and models as discussed in Sec.~\ref{sec:introduction}. In addition, MaskGIT~\citep{chang2022maskgit} adopts a heuristic training objective, which misses the $\frac{1}{t}$ term compared to Eq.~(\ref{eq:objective}), and lacks a theoretical link to maximum likelihood. We emphasize that it is precisely the theoretical foundation of maximum likelihood estimation that motivated us to scale discrete diffusion models for language modeling.



\subsection{Pre-training}
\label{sec:pre-traing}

LLaDA employs a Transformer~\citep{vaswani2017attention} as the mask predictor, similar to existing LLMs. However, LLaDA does not use a causal mask, as its formulation allows it to see the entire input for predictions.

We trained two variants of LLaDA with different sizes: 1B and 8B.
We summarize the model architecture of LLaDA 8B and LLaMA3 8B~\citep{dubey2024llama} here, and details are provided in Appendix~\ref{app:exp-1b-config}. We have ensured consistency in most hyperparameters while making several necessary modifications. We use vanilla multi-head attention instead of grouped query attention~\citep{ainslie2023gqa} for simplicity, as LLaDA is incompatible with KV caching, resulting in a different number of key and value heads. Consequently, the attention layer has more parameters, and we reduce the FFN dimension to maintain a comparable model size. Additionally, the vocabulary size differs due to a tokenizer~\citep{brown2020language} adapted on our data.

The LLaDA model is pre-trained on a dataset comprising 2.3 trillion (T) tokens, adhering to a data protocol that aligns closely with existing LLMs~\citep{qwen2,qwen2.5}, without the incorporation of any special techniques. The data are derived from online corpora, with low-quality content filtered through manually designed rules and LLM-based approaches. Beyond general text, the dataset encompasses high-quality code, math, and multilingual data. Please refer to Appendix~\ref{app:data} for more details about datasets. The mixing of data sources and domains is guided by scaled-down ARMs. The pre-training process utilizes a fixed sequence length of 4096 tokens, incurring a total computational cost of 0.13 million H800 GPU hours, similar to ARMs of the same scale and dataset size.
 

For a training sequence $x_0$, we randomly sample $t\in[0,1]$, mask each token independently with the same probability $t$ to obtain $x_t$ (see Fig.~\ref{fig:overview} (a)) and estimate Eq.~(\ref{eq:objective}) via the Monte Carlo method for stochastic gradient descent training. In addition, following~\citet{nie2024scaling}, to enhance the ability of LLaDA to handle variable-length data, we set 1\% of the pre-training data to a random length that is uniformly sampled from the range $[1, 4096]$.


We adopted the Warmup-Stable-Decay~\citep{hu2024minicpm} learning rate scheduler to monitor the training progress without interrupting continuous training. Specifically, we linearly increased the learning rate from 0 to \( 4 \times 10^{-4} \) over the first 2000 iterations and maintained it at \( 4 \times 10^{-4} \). After processing 1.2T tokens, we decayed the learning rate to \( 1 \times 10^{-4} \) and held it constant for the next 0.8T tokens to ensure stable training. Finally, we linearly reduced the learning rate from \( 1 \times 10^{-4} \) to \( 1 \times 10^{-5} \) for the last 0.3T tokens. Furthermore, we utilized the AdamW optimizer~\citep{loshchilov2017decoupled} with a weight decay of 0.1, a batch size of 1280, and a local batch size of $4$ per GPU. The 8B experiment was executed once, without any hyperparameter tuning.


\subsection{Supervised Fine-Tuning} 

We enhance the capability of LLaDA to follow instructions by supervised fine-tuning (SFT) with paired data \((p_0, r_0)\), where \(p_0\) is the prompt and \(r_0\) denotes the response. This is the simplest and most basic post-training method for LLMs. Technically, this requires to model the conditional distribution \(p_{\theta}(r_0|p_0)\) instead of \(p_{\theta}(x_0)\) in pre-training.

The implementation is similar to pre-training. As shown in Fig.~\ref{fig:overview} (b), we leave the prompt unchanged and mask the tokens in the response independently, as done for \(x_0\). Then, we feed both the prompt and the masked response \(r_t\) to the pre-trained mask predictor to compute the loss for SFT:
\begin{align}
\label{eq:sft-objective}
- \mathbb{E}_{t, p_0, r_0, r_t} \left[\frac{1}{t} \sum_{i=1}^{L'} \textbf{1}[r_t^i = \textrm{M}] \log p_{\theta}(r_0^i|p_0, r_t) \right],
\end{align}
where \(L'\) denotes a dynamic length specified later, and all other notations remain the same as before. 

Note that this approach is fully compatible with pre-training. Essentially, the concatenation of \(p_0\) and \(r_0\) can be treated as clean pre-training data \(x_0\), while the concatenation of \(p_0\) and \(r_t\) serves as the masked version \(x_t\). The process is identical to pre-training, with the only difference being that all masked tokens happen to appear in the \(r_0\) portion.

The LLaDA 8B model undergoes SFT on a dataset comprising 4.5 million pairs. Consistent with the pre-training process, both data preparation and training follow the SFT protocols utilized in existing LLMs~\citep{qwen2,qwen2.5}, without introducing any additional techniques to optimize LLaDA's performance.
The dataset spans multiple domains, including code, mathematics, and instruction-following. We append $|\text{EOS}|$ tokens to the end of short pairs in each mini-batch to ensure equal lengths across all data. We treat $|\text{EOS}|$ as a normal token during training and remove it during sampling, enabling LLaDA to control the response length automatically. Please refer to Appendix~\ref{app:data} for more details. 

We train for 3 epochs on the SFT data using a similar schedule to the pre-training phase. The learning rate is linearly increased from 0 to $2.5 \times 10^{-5}$ over the first 50 iterations and then kept constant. During the final $10\%$ of iterations, it is linearly reduced to $2.5 \times 10^{-6}$. Additionally, we set the weight decay to $0.1$, the global batch size to $256$, and the local batch size to $2$ per GPU. The SFT experiment was executed once, without any hyperparameter tuning.


\subsection{Inference}
\label{sec:inference}
As a generative model, LLaDA can sample new text and evaluate the likelihood of candidate text \emph{in a diffusion manner instead of the left-to-right autoregressive fashion}.

We begin with the reverse generation process. As illustrated in Fig.~\ref{fig:overview}~(c), given a prompt \(p_0\), we discretize the reverse process to sample from the model distribution \(p_\theta(r_0|p_0)\), starting from a fully masked response. The total number of sampling steps is a hyperparameter, which naturally provides LLaDA with a trade-off between efficiency and sample quality, as analyzed in Sec.~\ref{sec:analysis}. We employ uniformly distributed timesteps by default. In addition, the generation length is also treated as a hyperparameter, specifying the length of the fully masked sentence at the beginning of the sampling process. After generation, tokens appearing after the $|\text{EOS}|$ token are discarded. As detailed in Appendix~\ref{app:ablation-length}, since both pre-training and SFT are conducted using datasets with variable lengths, the final results are insensitive to this length hyperparameter.

At an intermediate step from time \(t \in (0, 1]\) to \(s \in [0, t)\), we feed both \(p_0\) and \(r_t\) into the mask predictor and predict all masked tokens simultaneously. Subsequently, we remask \(\frac{s}{t}\) of the predicted tokens in expectation to obtain \(r_s\), ensuring that the transition of the reverse process aligns with the forward process for accurate sampling~\citep{shi2024simplified, sahoo2024simple, ou2024your}. In principle, the remasking strategy should be purely random. However, inspired by the annealing tricks of sampling in LLMs~\citep{brown2020language,holtzman2019curious}, we adopt a low-confidence remasking strategy, where \(\frac{s}{t}\) of predicted tokens with the lowest confidence are remarked based on the predictions, same as the approach of~\citet{chang2022maskgit}.


% \cx{emphasize env. flexible sampling for free. after pretraining no fine-tuning xxx. diffusion sampling is better and it is the default method throughout the paper unless specified.}

We mention that LLaDA enables flexible sampling. In particular, it supports autoregressive and block diffusion~\citep{arriola2025block} sampling directly after the pre-training or SFT processes described above, without requiring any further modifications or training. We provide a detailed analysis in Appendix~\ref{app:sample}. Nevertheless, the diffusion sampling (i.e., the reverse generation process) yields the best performance and is adopted as the default throughout this paper, especially for all experiments presented in Sec.~\ref{sec:exp}.

For conditional likelihood evaluation, we can naturally utilize the upper bound in Eq.~(\ref{eq:sft-objective}). However, we find that the following equivalent form~\citep{ou2024your} exhibits lower variance and is more stable:
\begin{align}
\label{eq:ppl-eval}
    -\mathbb{E}_{l, r_0, r_l} \left[\frac{L}{l} \sum_{i=1}^L \textbf{1}[r_l^i = \textrm{M}] \log p_{\theta}(r_0^i|p_0, r_l) \right],
\end{align}
where $L$ is the sequence length of $r_0$, \( l \) is uniformly sampled from \( \{1, 2, \dots, L\} \), and \( r_l \) is obtained by uniformly sampling \( l \) tokens from \( r_0 \) without replacement for masking. 

We present the training and inference algorithms, along with theoretical details, in Appendix~\ref{app:formulation}.


\section{Experiments}
\label{sec:exp}

\begin{figure*}[t]
  \centering
  \begin{subfigure}{0.32\textwidth}
    \centering
    \includegraphics[width=\linewidth]{imgs/flops_mmlu_scatter.pdf}
  \end{subfigure}\hfill
  \begin{subfigure}{0.32\textwidth}
    \centering
    \includegraphics[width=\linewidth]{imgs/flops_arc_c_scatter.pdf}
  \end{subfigure}\hfill
  \begin{subfigure}{0.32\textwidth}
    \centering
    \includegraphics[width=\linewidth]{imgs/flops_cmmlu_scatter.pdf}
  \end{subfigure}
  \vspace{-0.2cm}
  \begin{subfigure}{0.32\textwidth}
    \centering
    \includegraphics[width=\linewidth]{imgs/flops_piqa_scatter.pdf}
  \end{subfigure}\hfill
  \begin{subfigure}{0.32\textwidth}
    \centering
    \includegraphics[width=\linewidth]{imgs/flops_GSM8K_scatter.pdf}
  \end{subfigure}\hfill
  \begin{subfigure}{0.32\textwidth}
    \centering
    \includegraphics[width=\linewidth]{imgs/flops_HumanEval_scatter.pdf}
  \end{subfigure}
  \vspace{0.1cm}
  \caption{\textbf{Scalability of LLaDA.} We evaluate the performance of LLaDA and our ARM baselines trained on the same data across increasing pre-training computational FLOPs. LLaDA exhibits strong scalability, matching the overall performance of ARMs on six tasks.}
  \label{fig:scaling}
  \vspace{-0.15cm}
\end{figure*}


We evaluate the scalability, instruction-following, and in-context learning capabilities of LLaDA on standard benchmarks, followed by analyses and case studies to provide a comprehensive assessment.

\subsection{Scalability of LLaDA on Language Tasks}
\label{sec:scaling}

We first investigate the \emph{scalability} of LLaDA on downstream tasks in comparison with the ARM baselines we constructed. Specifically, at the 1B scale, we ensured that LLaDA and ARM shared the same architecture, data, and all other configurations. At larger scales, we also report results for LLaDA and ARM models of slightly different sizes trained on the same data due to resource limitations. Please refer to Appendix~\ref{app:exp-1b-config} for more details. We use the pre-training computational cost as a unified scaling metric. For evaluation, we focused on six standard and diverse tasks.

Fig.~\ref{fig:scaling} shows that LLaDA demonstrates impressive scalability, with its overall trend highly competitive with ARMs. Notably, on tasks such as MMLU and GSM8K, LLaDA exhibits even stronger scalability. Even on relatively weaker tasks like PIQA, the performance gap with ARMs narrows as scale increases. To account for the influence of outliers, we opted not to fit quantitative curves, avoiding potential misinterpretation. Nevertheless, the results clearly demonstrate the scalability of LLaDA. 

Considering LLaDA’s advantages on certain benchmarks, we hypothesize that this performance gain stems from a key architectural difference: while autoregressive models optimize only left-to-right conditional probabilities, LLaDA is trained to consider multiple conditioning directions, as detailed in Appendix~\ref{app:for-infer}, which may offer greater flexibility and lead to better generalization. This hypothesis is motivated by LLaDA’s strong performance on reversal reasoning in Sec.~\ref{sec:analysis} and the ablation studies on sampling strategies in Appendix~\ref{app:sample}.


\citet{nie2024scaling} suggests that MDM requires 16 times more computation than ARM to achieve the same likelihood. However, key differences make our findings more broadly applicable. In particular, likelihood is a relatively indirect metric for downstream task performance, and diffusion optimizes a bound of the likelihood, making it not directly comparable to ARM. Additionally, we extended the scaling range from $10^{18}\sim10^{20}$ FLOPs in~\citet{nie2024scaling} to $10^{20} \sim 10^{23}$ FLOPs in this work.




\subsection{Benchmark Results}
To comprehensively evaluate the \emph{in-context learning} and \emph{instruction-following} capabilities of LLaDA 8B, we conducted detailed comparisons with existing LLMs~\citep{dubey2024llama,touvron2023llama2, qwen2,qwen2.5, bi2024deepseek,jiang2023mistral} of similar scale. Task selection and evaluation protocols followed existing studies, covering popular benchmarks in general tasks, mathematics, code, and Chinese. Further details are provided in Appendix~\ref{app:bench}. For a more direct comparison, we re-evaluated representative LLMs~\citep{dubey2024llama, touvron2023llama2} in our implementation.

As shown in Tab.~\ref{tab:base}, after pretraining on 2.3T tokens, LLaDA 8B Base demonstrates remarkable performance, surpassing LLaMA2 7B Base on nearly all tasks, and is overall competitive with LLaMA3 8B Base. LLaDA shows advantages in math and Chinese tasks. We conjecture that the strengths stem from the same factors as its relatively weaker performance in some tasks—differences in data quality and distribution, largely due to the closed-source situation of LLM datasets. 


\begin{table*}[t!]
    \centering
    \caption{\textbf{Benchmark Results of Pre-trained LLMs.} $^{*}$ indicates that models are evaluated under the same protocol, detailed in Appendix~\ref{app:bench}. Results indicated by $^{\dagger}$ and  $^{\mathparagraph}$ are sourced from~\citet{qwen2, qwen2.5} and~\citet{bi2024deepseek} respectively. The numbers in parentheses represent the number of shots used for in-context learning. ``-'' indicates unknown data.}
    \label{tab:base}
    \vspace{.2cm} 
    \begin{adjustbox}{max width=\textwidth}
    \begin{tabular}{l|ccc|cccc}
      \toprule
         & LLaDA 8B$^{*}$ & LLaMA3 8B$^{*}$ & LLaMA2 7B$^{*}$ &  Qwen2 7B$^{\dagger}$ & Qwen2.5 7B$^{\dagger}$ & Mistral 7B$^{\dagger}$ & Deepseek 7B$^{\mathparagraph}$  \\
      \midrule
      Model & Diffusion & AR & AR & AR & AR & AR & AR \\
      Training tokens & 2.3T & 15T & 2T & 7T & 18T & - & 2T \\
      \midrule
         \multicolumn{8}{c}{General Tasks}\\
      \midrule
          MMLU & \textbf{65.9} (5) & 65.4 (5) & 45.9 (5) & 70.3 (5) & 74.2 (5) & 64.2 (5) & 48.2 (5) \\
          BBH & 49.7 (3) & \textbf{62.1} (3) & 39.4 (3) & 62.3 (3) & 70.4 (3) & 56.1 (3) & 39.5 (3) \\
          ARC-C & 45.9 (0) & \textbf{53.1} (0) & 46.3 (0) & 60.6 (25) & 63.7 (25) & 60.0 (25) & 48.1 (0) \\
          Hellaswag & 70.5 (0) & \textbf{79.1} (0) & 76.0 (0) & 80.7 (10) & 80.2 (10) & 83.3 (10) & 75.4 (0) \\
          TruthfulQA & \textbf{46.1} (0) & 44.0 (0) & 39.0 (0) & 54.2 (0) & 56.4 (0) & 42.2 (0) & - \\ 
          WinoGrande & 74.8 (5) & \textbf{77.3} (5) & 72.5 (5) & 77.0 (5) & 75.9 (5) & 78.4 (5) & 70.5 (0) \\
          PIQA & 73.6 (0) & \textbf{80.6} (0) & 79.1 (0) & - & - & - & 79.2 (0) \\
      \midrule
        \multicolumn{8}{c}{Mathematics  \& Science}\\
      \midrule
        GSM8K & \textbf{70.3} (4) & 48.7 (4) & 13.1 (4) & 80.2 (4) & 85.4 (4) & 36.2 (4) & 17.4 (8) \\
        Math & \textbf{31.4} (4) & 16.0 (4) & 4.3 (4) & 43.5 (4) & 49.8 (4) & 10.2 (4) & 6.0 (4) \\
        GPQA & 25.2 (5) & \textbf{25.9} (5) & 25.7 (5) & 30.8 (5) & 36.4 (5) & 24.7 (5) & - \\
      \midrule
        \multicolumn{8}{c}{Code} \\
      \midrule
        HumanEval & \textbf{35.4} (0) & 34.8 (0) & 12.8 (0) & 51.2 (0) & 57.9 (0) & 29.3 (0) & 26.2 (0) \\    
        HumanEval-FIM & \textbf{73.8} (2) & 73.3 (2) & 26.9 (2) & - & - & - & - \\
        MBPP & 40.0 (4) & \textbf{48.8} (4) & 23.2 (4) & 64.2 (0) & 74.9 (0) & 51.1 (0) & 39.0 (3) \\
      \midrule
        \multicolumn{8}{c}{Chinese}\\
      \midrule
        CMMLU & \textbf{69.9} (5) & 50.7 (5) & 32.5 (5) & 83.9 (5) & - & - & 47.2 (5) \\
        C-Eval & \textbf{70.5} (5) & 51.7 (5) & 34.0 (5) & 83.2 (5) & - & - & 45.0 (5) \\ 
      \bottomrule
    \end{tabular}
    \end{adjustbox}
\end{table*}


 \begin{table*}[t!]
    \centering
    \caption{\textbf{Benchmark Results of Post-trained LLMs.} LLaDA only employs an SFT procedure, while other models have extra reinforcement learning (RL) alignment.  $^{*}$ indicates models are evaluated under the same protocol, detailed in Appendix~\ref{app:bench}. Results indicated by $^{\dagger}$ and  $^{\mathparagraph}$ are sourced from~\citet{qwen2.5} and~\citet{bi2024deepseek} respectively. The numbers in parentheses represent the number of shots used for in-context learning. ``-'' indicates unknown data.}
    \label{tab:chat}
    \vspace{.2cm}
    \begin{adjustbox}{max width=\textwidth}
    \begin{tabular}{l|ccc|cccc}
      \toprule
  & LLaDA 8B$^{*}$ & LLaMA3 8B$^{*}$ & LLaMA2 7B$^{*}$ &  Qwen2 7B$^{\dagger}$ & Qwen2.5 7B$^{\dagger}$ & Gemma2 9B$^{\dagger}$ & Deepseek 7B$^{\mathparagraph}$  \\
      \midrule
Model & Diffusion & AR & AR & AR & AR & AR & AR \\
      Training tokens & 2.3T & 15T & 2T & 7T & 18T & 8T & 2T \\
      \midrule
      Post-training & SFT & SFT+RL & SFT+RL &SFT+RL &SFT+RL &SFT+RL &SFT+RL \\
      Alignment pairs & 4.5M & - & - & 0.5M + - & 1M + 0.15M & - & 1.5M + -\\
      \midrule
         \multicolumn{8}{c}{General Tasks}\\
      \midrule
          MMLU & 65.5 (5)  & \textbf{68.4} (5) & 44.1 (5) & - & - & - & 49.4 (0) \\
          MMLU-pro & 37.0 (0) & \textbf{41.9} (0) & 4.6 (0) & 44.1 (5) & 56.3 (5) & 52.1 (5) & - \\
          Hellaswag & 74.6 (0) & \textbf{75.5} (0) & 51.5 (0) & - & - & - & 68.5 (-)\\
       ARC-C & \textbf{88.5} (0) & 82.4 (0) & 57.3 (0) & - & - & - & 49.4 (-) \\
      \midrule
        \multicolumn{8}{c}{Mathematics \& Science}\\
      \midrule
        GSM8K & 69.4 (4) & \textbf{78.3} (4) & 29.0 (4) & 85.7 (0) & 91.6 (0) & 76.7 (0) & 63.0 (0)\\
        Math & \textbf{31.9} (0) & 29.6 (0) & 3.8 (0) & 52.9 (0) & 75.5 (0) & 44.3 (0) & 15.8 (0)\\
        GPQA & \textbf{33.3} (5) & 31.9 (5) & 28.4 (5) & 34.3 (0) & 36.4 (0) & 32.8 (0) & -\\
      \midrule
        \multicolumn{8}{c}{Code} \\
      \midrule
        HumanEval & 49.4 (0) & \textbf{59.8} (0) & 16.5 (0) & 79.9 (0) & 84.8 (0) & 68.9 (0) & 48.2 (-) \\     
        MBPP & 41.0 (4) & \textbf{57.6} (4) & 20.6 (4) & 67.2 (0) & 79.2 (0) & 74.9 (0) & 35.2 (-)\\
      \bottomrule
    \end{tabular}
    \end{adjustbox}
\end{table*}

Notably, we have carefully ruled out the possibility of data leakage by taking GSM8K as an example. First, as shown in Fig.~\ref{fig:scaling}, LLaDA outperformed ARM baselines regarding GSM8K. Moreover, the conclusion remains on a fully unseen GSM8K-like task~\citep{YXLA2024-gsm1} in Appendix~\ref{app:igsm}.

Further, Tab.~\ref{tab:chat} compares the performance of LLaDA 8B Instruct with existing LLMs. SFT improved LLaDA's performance on most downstream tasks. A few metrics, such as MMLU, showed declines, possibly due to the suboptimal quality of the SFT data. Overall, since we did not perform alignment with reinforcement learning (RL), our results are slightly behind LLaMA3 8B Instruct, though the gaps in many metrics remain small. Notably, even with only SFT, LLaDA demonstrates impressive instruction-following abilities, as detailed in Sec.~\ref{sec:case}. We leave RL-based alignment for future work.

All results in Sec.~\ref{sec:exp} are based on pure diffusion methods, as they achieve better overall performance than approaches incorporating autoregressive components. Specifically, we use Eq.~(\ref{eq:ppl-eval}) for conditional likelihood estimation and apply low-confidence remasking for sampling. For LLaDA 8B Instruct, block diffusion style sampling performs better on GSM8K and Math, with scores of 78.6 and 42.2, compared to 69.4 and 31.9 in Tab.~\ref{tab:chat}. This gain is due to extensive $|\text{EOS}|$ token padding in the SFT data, causing early termination in low-confidence remasking. Please refer to Appendix~\ref{app:sample} for details.

Overall, despite the lack of data transparency, we have made every effort to adopt standardized procedures and introduce diverse tasks, we believe they sufficiently demonstrate the extraordinary capabilities of LLaDA, which is the only competitive non-autoregressive model to our knowledge.



\subsection{Reversal Reasoning and Analyses}
\label{sec:analysis}  
\begin{table*}[t!]
\centering
\caption{\textbf{Visualization of the Sampling Process and a Generated Multi-round Dialogue.} In the response of LLaDA, darker colors indicate tokens predicted in the later stages of sampling, while lighter colors correspond to earlier predictions.}
\label{tab:visualization}
\vspace{5pt} % 调整这里的数值来控制距离
\scalebox{1}{
\begin{adjustbox}{max width=\textwidth}
\begin{tabular}{l p{12.5cm} }
\toprule
\emph{Sampling Process} & \\
\midrule
~~~~User & Lily can run 12 kilometers per hour for 4 hours. After that, she runs 6 kilometers per hour. How many kilometers can she run in 8 hours?
\\
~~~~LLaDA & \includegraphics[width=.85\textwidth,valign=t]{imgs/visualization/diff_math.png}\\
\midrule
 \multicolumn{2}{l}{\emph{Multi-round Dialogue Case}}  \\
\midrule
~~~~User & Could you please share with me what the initial two lines of the renowned poem `The Road Not Taken'?\\
~~~~LLaDA & Certainly! The initial two lines of the renowned poem ``The Road Not Taken" by Robert Frost are: ``Two roads diverged in a yellow wood, And sorry I could not travel both." These lines set the stage for the poem, introducing the idea of a speaker facing a choice between two paths. \\
~~~~User & Please help me translate into Chinese. \\
~~~~LLaDA & Sure! The first two lines of ``The Road Not Taken" by Robert Frost can be translated into Chinese as: \begin{CJK}{UTF8}{gbsn}``两条路分岔在黄色的树林中，遗憾我不能同时走''\end{CJK}\\
~~~~User & Please translate into German. \\
~~~~LLaDA & Sure! The first two lines of ``The Road Not Taken" by Robert Frost can be translated into German as: ``Zwei Wege trennten sich im gelben Wald, und leider konnte ich nicht beide tun." \\
~~~~User & Please also write a poem about life choices. I require 5 sentences, and the first word of each sentence starts with C as the first letter. \\
~~~~LLaDA & Certainly! Here's a poem about life choices, with each sentence starting with the letter ``C": Choices cascade through life's journey, Creating connections with each decision, Careful consideration leads to better paths, Courageous choices shape our destiny, Choosing wisely guides us through the years. \\
\bottomrule
\end{tabular}
\end{adjustbox}
}
\end{table*}

To quantify the reversal reasoning~\citep{berglund2023reversal} ability of models, we follow the protocol established in \citet{AL2023-knowledge2}. Specifically, we construct a dataset of 496 famous Chinese poem sentence pairs. Given a sentence from a poem, models are tasked with generating the subsequent line (forward) or the preceding line (reversal) without additional fine-tuning. Examples can be found in~\cref{app:exp-poem}. This setting provides a straightforward and more realistic evaluation compared to previous studies~\citep{nie2024scaling,kitouni2024factorization}.

As shown in Tab.~\ref{tbl:poems_completion}, LLaDA effectively addresses the \emph{reversal curse}~\citep{berglund2023reversal}, demonstrating consistent zero-shot performance across both forward and reversal tasks. In contrast, both Qwen 2.5 and GPT-4o exhibit a significant gap between the two. The results on forward generation confirm that both ARMs are strong, benefiting from significantly larger datasets and greater computational resources than LLaDA. However, LLaDA outperforms both by a large margin in the reversal task.

\begin{wraptable}{r}{0.47\textwidth}  % 表格靠右，占 45% 宽
  \centering
  \caption{\textbf{Comparison on the Poem Completion task.} }
  \label{tbl:poems_completion}
  \begin{adjustbox}{max width=\linewidth}
  \begin{tabular}{lcc}
    \toprule
    & Forward & Reversal \\ 
    \midrule
    GPT-4o (2024‑08‑06)  & \textbf{82.7} & 34.3 \\
    Qwen2.5‑7B Instruct & 75.9          & 38.0 \\
    LLaDA‑8B Instruct   & 51.8          & \textbf{45.6} \\
    \bottomrule
  \end{tabular}
  \end{adjustbox}
\end{wraptable}



We did not design anything special for reversal tasks. Intuitively, LLaDA treats tokens uniformly without inductive bias, leading to balanced performance. See Appendix~\ref{app:for-infer} for details.



We also analyze the effect of different sampling strategies for LLaDA, including autoregressive sampling, block diffusion~\citep{arriola2025block} sampling, and pure diffusion sampling, showing that pure diffusion sampling achieves the best overall performance, as detailed in Appendix~\ref{app:sample}. 

In addition, we examine LLaDA’s sampling speed and memory consumption, showing that it enables a flexible trade-off between generation quality and speed. See Appendix~\ref{app:analysis} for more details.

Classifier-free guidance (CFG)~\citep{ho2022classifier,nie2024scaling} is a widely used technique in diffusion models to improve generation quality. To ensure a fair comparison with ARMs, we do not apply CFG to LLaDA in the main text. However, we show that LLaDA is compatible with CFG and consistently benefits from its application. See Appendix~\ref{app:cfg} for more details.


\subsection{Case Studies}
\label{sec:case}

We present samples generated by LLaDA 8B Instruct in Tab.~\ref{tab:visualization}, showcasing its instruction-following capabilities. First, the table illustrates LLaDA’s ability to generate coherent, fluent, and extended text in a non-autoregressive manner. Second, it highlights the model’s multi-turn dialogue capability, effectively retaining conversation history and producing contextually appropriate responses across multiple languages. Such \emph{chat} capabilities of LLaDA are impressive, as it departs from conventional ARMs for the first time, to the best of our knowledge.
See more case studies in Appendix~\ref{app:more_case}. 


\section{Related Work} 
Diffusion models~\citep{sohl2015deep,ho2020denoising,song2020score} have achieved remarkable success in visual domains but remain unverified for large-scale (e.g., models trained with over $10^{23}$ FLOPs)  language modeling, despite growing interest and extensive research efforts.

A simple approach is to continuousize text data and apply continuous diffusion models directly~\citep{li2022diffusion,gong2022diffuseq,han2022ssd,strudel2022self,chen2022analog,dieleman2022continuous,richemond2022categorical,wu2023ardiffusion,mahabadi2024tess,ye2023dinoiser,zhang2023planner}. Alternatively, some methods model continuous parameters of discrete distributions instead~\citep{lou2023reflected,graves2023bayesian,lin2023text,xue2024unifying, zhang2025target}. However, scalability remains a significant challenge for these approaches. For instance, a 1B model may require 64 times the compute of an ARM to achieve comparable performance~\citep{gulrajani2024likelihood}.



Another approach replaces continuous diffusion with discrete processes featuring new forward and reverse dynamics, leading to numerous variants~\citep{hoogeboom2021argmax,hoogeboom2021autoregressive,he2022diffusionbert,campbell2022continuous,meng2022concrete,reid2022diffuser,sun2022score,kitouni2023disk,Zheng2023ARD,chen2023fast,ye2023diffusion,gat2024discrete,zheng2024maskeddiffusionmodelssecretly, kapur2024diffusion}. The original diffusion model paper~\citep{sohl2015deep} introduced both continuous-state and discrete-state transition kernels under a unified diffusion framework. \citet{austin2021structured} was among the pioneering works that introduced discrete diffusion models into language modeling, demonstrating the feasibility of this approach. \citet{lou2023discrete} showed that masked diffusion, as a special case of discrete diffusion, achieves perplexity comparable to or surpassing ARMs at GPT-2 scale. \citet{shi2024simplified, sahoo2024simple, ou2024your} established fundamental theoretical results, which motivated our model design, training, and inference (see Appendix~\ref{app:formulation} for details).  \citet{nie2024scaling} introduced the scaling laws for MDMs in language modeling and explored how MDMs can be leveraged for language tasks such as question answering at the GPT-2 scale. \citet{gong2024scaling} demonstrated the potential of fine-tuning an ARM within the MDM framework. However, the improvements observed by~\citet{gong2024scaling} are limited to specific metrics, and their approach does not address the performance achievable through pure diffusion-based training. Concurrent work~\citep{khanna2025mercury} demonstrates the potential of diffusion language models in code generation and highlights their advantages in inference efficiency. Nonetheless, as it is a closed-source product, specific details such as training procedures and sampling methods remain unknown.

In comparison, this study scales MDM to an unprecedented size of 8B parameters from scratch, achieving performance comparable to leading LLMs such as LLaMA 3.

Additionally, a parallel line of work on image generation~\citep{chang2022maskgit, chang2023muse, you2025effective} aligns well with the application of MDMs to text data. Moreover, MDMs have also shown promise in other domains such as protein generation~\citep{wang2024diffusion, wang2024dplm}, where they have achieved promising results. Notably, a series of studies~\citep{arriola2025block, kou2024cllms, xu2025show, liu2024think, zhu2025di, ren2025fast, hayakawa2024distillation, zhao2024informed, zheng2024masked, park2024jump, deschenaux2024beyond} have explored techniques such as architectural optimization, distillation, and sampling algorithm design to accelerate MDMs sampling.

\section{Conclusion and Discussion}
\label{sec:conclusion}
We introduce LLaDA, a diffusion language model trained from scratch with an unprecedented scale of 8B parameters. LLaDA demonstrates strong capabilities in scalability,  in-context learning, and instruction-following, achieving performance comparable to strong LLMs such as LLaMA3. In addition, LLaDA offers unique advantages, such as bidirectional modeling and enhanced robustness, effectively addressing the relevant limitations of existing LLMs. Our findings show the promise of diffusion models for language modeling at scale and challenge the common assumption that these essential capabilities are inherently tied to ARMs. These results represent a new paradigm for language modeling and uncover novel insights, demonstrating a high degree of scientific innovation.

\textbf{Limitations.} While promising, the full potential of diffusion models remains to be fully explored. Several limitations of this work present significant opportunities for future research. The generation length is a user-specified hyperparameter. Although LLaDA is insensitive to this hyperparameter as detailed in Appendix~\ref{app:ablation-length}, we believe that adopting an adaptive generation length would offer a more efficient solution. Due to computational constraints, direct comparisons between LLaDA and ARMs---such as training on identical datasets---were restricted to a computational budget of less than $10^{23}$ FLOPs. To allocate resources for training the largest possible LLaDA model and showcasing its potential, we were unable to scale the ARM baseline to the same extent. Moreover, no specialized attention mechanisms or position embeddings were designed for LLaDA, nor were any system-level architectural optimizations such as KV cache applied. On the inference side, more efficient and controllable~\citep{ho2022classifier, dhariwal2021diffusion, schiff2024simple} sampling algorithms remain preliminary. Furthermore, LLaDA has yet to undergo alignment with reinforcement learning~\citep{ouyang2022training,rafailov2024direct}, which is crucial for improving its performance and alignment with human intent.

Looking ahead, both the model scale and the amount of training data for LLaDA remain smaller than those of leading ARM counterparts~\citep{dubey2024llama,qwen2.5, achiam2023gpt,gemini1_5,claude35sonnet,liu2024deepseek}, highlighting the need for further scaling to fully evaluate its capabilities. In addition, LLaDA's ability to process multi-modal data remains unexplored. Its impact on prompt tuning techniques~\citep{wei2022chain} and integration into agent-based systems~\citep{park2023generative,wang2024survey} is still not fully understood. Finally, a systematic investigation into post-training for LLaDA (e.g., O1-like systems~\citep{o1,guo2025deepseek}) is needed to further unlock the potential of diffusion language models.


\section*{Acknowledgements}
This work was supported by the National Natural Science Foundation of China (No. 92470118); Beijing Natural Science Foundation (No. L247030); Beijing Nova Program (No. 20220484044); and Ant Group Research Fund.



%%%%%%%%%%%%%%%%%%%%%%%%%%%%%%%%%%%%%%%%%%%%%%%%%%%%%%%%%%%%

\bibliographystyle{unsrtnat}
\bibliography{references}







%%%%%%%%%%%%%%%%%%%%%%%%%%%%%%%%%%%%%%%%%%%%%%%%%%%%%%%%%%%
%%%%%%%%%%%%%%%%%%%%%%%%%%%%%%%%%%%%%%%%%%%%%%%%%%%%%%%%%%%%
% \section*{NeurIPS Paper Checklist}

% \begin{enumerate}

% \item {\bf Claims}
%     \item[] Question: Do the main claims made in the abstract and introduction accurately reflect the paper's contributions and scope?
%     \item[] Answer: \answerYes{} % Replace by \answerYes{}, \answerNo{}, or \answerNA{}.
%     \item[] Justification: The contributions and scope of the paper are well summarized in the abstract and introduction.
%     \item[] Guidelines:
%     \begin{itemize}
%         \item The answer NA means that the abstract and introduction do not include the claims made in the paper.
%         \item The abstract and/or introduction should clearly state the claims made, including the contributions made in the paper and important assumptions and limitations. A No or NA answer to this question will not be perceived well by the reviewers. 
%         \item The claims made should match theoretical and experimental results, and reflect how much the results can be expected to generalize to other settings. 
%         \item It is fine to include aspirational goals as motivation as long as it is clear that these goals are not attained by the paper. 
%     \end{itemize}

% \item {\bf Limitations}
%     \item[] Question: Does the paper discuss the limitations of the work performed by the authors?
%     \item[] Answer: \answerYes{} % Replace by \answerYes{}, \answerNo{}, or \answerNA{}.
%     \item[] Justification: We discuss the limitations of our work in Sec.~\ref{sec:conclusion}.
%     \item[] Guidelines:
%     \begin{itemize}
%         \item The answer NA means that the paper has no limitation while the answer No means that the paper has limitations, but those are not discussed in the paper. 
%         \item The authors are encouraged to create a separate "Limitations" section in their paper.
%         \item The paper should point out any strong assumptions and how robust the results are to violations of these assumptions (e.g., independence assumptions, noiseless settings, model well-specification, asymptotic approximations only holding locally). The authors should reflect on how these assumptions might be violated in practice and what the implications would be.
%         \item The authors should reflect on the scope of the claims made, e.g., if the approach was only tested on a few datasets or with a few runs. In general, empirical results often depend on implicit assumptions, which should be articulated.
%         \item The authors should reflect on the factors that influence the performance of the approach. For example, a facial recognition algorithm may perform poorly when image resolution is low or images are taken in low lighting. Or a speech-to-text system might not be used reliably to provide closed captions for online lectures because it fails to handle technical jargon.
%         \item The authors should discuss the computational efficiency of the proposed algorithms and how they scale with dataset size.
%         \item If applicable, the authors should discuss possible limitations of their approach to address problems of privacy and fairness.
%         \item While the authors might fear that complete honesty about limitations might be used by reviewers as grounds for rejection, a worse outcome might be that reviewers discover limitations that aren't acknowledged in the paper. The authors should use their best judgment and recognize that individual actions in favor of transparency play an important role in developing norms that preserve the integrity of the community. Reviewers will be specifically instructed to not penalize honesty concerning limitations.
%     \end{itemize}

% \item {\bf Theory assumptions and proofs}
%     \item[] Question: For each theoretical result, does the paper provide the full set of assumptions and a complete (and correct) proof?
%     \item[] Answer: \answerYes{} % Replace by \answerYes{}, \answerNo{}, or \answerNA{}.
%     \item[] Justification: In Appendix~\ref{app:formulation}, we provide an overview of the theoretical foundations required for this work, along with precise citations. Readers may refer to our overview or consult the original references for full theoretical details.
%     \item[] Guidelines:
%     \begin{itemize}
%         \item The answer NA means that the paper does not include theoretical results. 
%         \item All the theorems, formulas, and proofs in the paper should be numbered and cross-referenced.
%         \item All assumptions should be clearly stated or referenced in the statement of any theorems.
%         \item The proofs can either appear in the main paper or the supplemental material, but if they appear in the supplemental material, the authors are encouraged to provide a short proof sketch to provide intuition. 
%         \item Inversely, any informal proof provided in the core of the paper should be complemented by formal proofs provided in appendix or supplemental material.
%         \item Theorems and Lemmas that the proof relies upon should be properly referenced. 
%     \end{itemize}

%     \item {\bf Experimental result reproducibility}
%     \item[] Question: Does the paper fully disclose all the information needed to reproduce the main experimental results of the paper to the extent that it affects the main claims and/or conclusions of the paper (regardless of whether the code and data are provided or not)?
%     \item[] Answer: \answerYes{} % Replace by \answerYes{}, \answerNo{}, or \answerNA{}.
%     \item[] Justification: We provide details on data collection, training setup, model architecture, and evaluation protocols in Sec.~\ref{sec:approach} and Appendix~\ref{app:experiment}.
%     \item[] Guidelines:
%     \begin{itemize}
%         \item The answer NA means that the paper does not include experiments.
%         \item If the paper includes experiments, a No answer to this question will not be perceived well by the reviewers: Making the paper reproducible is important, regardless of whether the code and data are provided or not.
%         \item If the contribution is a dataset and/or model, the authors should describe the steps taken to make their results reproducible or verifiable. 
%         \item Depending on the contribution, reproducibility can be accomplished in various ways. For example, if the contribution is a novel architecture, describing the architecture fully might suffice, or if the contribution is a specific model and empirical evaluation, it may be necessary to either make it possible for others to replicate the model with the same dataset, or provide access to the model. In general. releasing code and data is often one good way to accomplish this, but reproducibility can also be provided via detailed instructions for how to replicate the results, access to a hosted model (e.g., in the case of a large language model), releasing of a model checkpoint, or other means that are appropriate to the research performed.
%         \item While NeurIPS does not require releasing code, the conference does require all submissions to provide some reasonable avenue for reproducibility, which may depend on the nature of the contribution. For example
%         \begin{enumerate}
%             \item If the contribution is primarily a new algorithm, the paper should make it clear how to reproduce that algorithm.
%             \item If the contribution is primarily a new model architecture, the paper should describe the architecture clearly and fully.
%             \item If the contribution is a new model (e.g., a large language model), then there should either be a way to access this model for reproducing the results or a way to reproduce the model (e.g., with an open-source dataset or instructions for how to construct the dataset).
%             \item We recognize that reproducibility may be tricky in some cases, in which case authors are welcome to describe the particular way they provide for reproducibility. In the case of closed-source models, it may be that access to the model is limited in some way (e.g., to registered users), but it should be possible for other researchers to have some path to reproducing or verifying the results.
%         \end{enumerate}
%     \end{itemize}


% \item {\bf Open access to data and code}
%     \item[] Question: Does the paper provide open access to the data and code, with sufficient instructions to faithfully reproduce the main experimental results, as described in supplemental material?
%     \item[] Answer: \answerNo{} % Replace by \answerYes{}, \answerNo{}, or \answerNA{}.
%     \item[] Justification: As with most large language models, we do not provide open access to the training data at this stage due to data ownership restrictions. We provide details on data collection, training setup, model architecture, and evaluation protocols in Sec.~\ref{sec:approach} and Appendix~\ref{app:experiment}. In addition, we will release the model weights and evaluation code upon acceptance.
%     \item[] Guidelines:
%     \begin{itemize}
%         \item The answer NA means that paper does not include experiments requiring code.
%         \item Please see the NeurIPS code and data submission guidelines (\url{https://nips.cc/public/guides/CodeSubmissionPolicy}) for more details.
%         \item While we encourage the release of code and data, we understand that this might not be possible, so “No” is an acceptable answer. Papers cannot be rejected simply for not including code, unless this is central to the contribution (e.g., for a new open-source benchmark).
%         \item The instructions should contain the exact command and environment needed to run to reproduce the results. See the NeurIPS code and data submission guidelines (\url{https://nips.cc/public/guides/CodeSubmissionPolicy}) for more details.
%         \item The authors should provide instructions on data access and preparation, including how to access the raw data, preprocessed data, intermediate data, and generated data, etc.
%         \item The authors should provide scripts to reproduce all experimental results for the new proposed method and baselines. If only a subset of experiments are reproducible, they should state which ones are omitted from the script and why.
%         \item At submission time, to preserve anonymity, the authors should release anonymized versions (if applicable).
%         \item Providing as much information as possible in supplemental material (appended to the paper) is recommended, but including URLs to data and code is permitted.
%     \end{itemize}


% \item {\bf Experimental setting/details}
%     \item[] Question: Does the paper specify all the training and test details (e.g., data splits, hyperparameters, how they were chosen, type of optimizer, etc.) necessary to understand the results?
%     \item[] Answer: \answerYes{} % Replace by \answerYes{}, \answerNo{}, or \answerNA{}.
%     \item[] Justification: We provide details on data collection, training setup, model architecture, and evaluation protocols in Sec.~\ref{sec:approach} and Appendix~\ref{app:experiment}. 
%     \item[] Guidelines:
%     \begin{itemize}
%         \item The answer NA means that the paper does not include experiments.
%         \item The experimental setting should be presented in the core of the paper to a level of detail that is necessary to appreciate the results and make sense of them.
%         \item The full details can be provided either with the code, in appendix, or as supplemental material.
%     \end{itemize}

% \item {\bf Experiment statistical significance}
%     \item[] Question: Does the paper report error bars suitably and correctly defined or other appropriate information about the statistical significance of the experiments?
%     \item[] Answer: \answerNo{} % Replace by \answerYes{}, \answerNo{}, or \answerNA{}.
%     \item[] Justification: Due to the high cost of model training, it is difficult for us to run repeated experiments.
%     \item[] Guidelines:
%     \begin{itemize}
%         \item The answer NA means that the paper does not include experiments.
%         \item The authors should answer "Yes" if the results are accompanied by error bars, confidence intervals, or statistical significance tests, at least for the experiments that support the main claims of the paper.
%         \item The factors of variability that the error bars are capturing should be clearly stated (for example, train/test split, initialization, random drawing of some parameter, or overall run with given experimental conditions).
%         \item The method for calculating the error bars should be explained (closed form formula, call to a library function, bootstrap, etc.)
%         \item The assumptions made should be given (e.g., Normally distributed errors).
%         \item It should be clear whether the error bar is the standard deviation or the standard error of the mean.
%         \item It is OK to report 1-sigma error bars, but one should state it. The authors should preferably report a 2-sigma error bar than state that they have a 96\% CI, if the hypothesis of Normality of errors is not verified.
%         \item For asymmetric distributions, the authors should be careful not to show in tables or figures symmetric error bars that would yield results that are out of range (e.g. negative error rates).
%         \item If error bars are reported in tables or plots, The authors should explain in the text how they were calculated and reference the corresponding figures or tables in the text.
%     \end{itemize}

% \item {\bf Experiments compute resources}
%     \item[] Question: For each experiment, does the paper provide sufficient information on the computer resources (type of compute workers, memory, time of execution) needed to reproduce the experiments?
%     \item[] Answer: \answerYes{} % Replace by \answerYes{}, \answerNo{}, or \answerNA{}.
%     \item[] Justification: We provide the training computer resources in Sec.~\ref{sec:approach}. We also provide an inference efficiency analysis in Appendix~\ref{app:analysis}. 
%     \item[] Guidelines:
%     \begin{itemize}
%         \item The answer NA means that the paper does not include experiments.
%         \item The paper should indicate the type of compute workers CPU or GPU, internal cluster, or cloud provider, including relevant memory and storage.
%         \item The paper should provide the amount of compute required for each of the individual experimental runs as well as estimate the total compute. 
%         \item The paper should disclose whether the full research project required more compute than the experiments reported in the paper (e.g., preliminary or failed experiments that didn't make it into the paper). 
%     \end{itemize}
    
% \item {\bf Code of ethics}
%     \item[] Question: Does the research conducted in the paper conform, in every respect, with the NeurIPS Code of Ethics \url{https://neurips.cc/public/EthicsGuidelines}?
%     \item[] Answer: \answerYes{} % Replace by \answerYes{}, \answerNo{}, or \answerNA{}.
%     \item[] Justification: The authors have reviewed and followed the NeurIPS Code of Ethics.
%     \item[] Guidelines:
%     \begin{itemize}
%         \item The answer NA means that the authors have not reviewed the NeurIPS Code of Ethics.
%         \item If the authors answer No, they should explain the special circumstances that require a deviation from the Code of Ethics.
%         \item The authors should make sure to preserve anonymity (e.g., if there is a special consideration due to laws or regulations in their jurisdiction).
%     \end{itemize}


% \item {\bf Broader impacts}
%     \item[] Question: Does the paper discuss both potential positive societal impacts and negative societal impacts of the work performed?
%     \item[] Answer: \answerYes{} % Replace by \answerYes{}, \answerNo{}, or \answerNA{}.
%     \item[] Justification: We provide a discussion about societal impacts in the Appendix~\ref{app:impact}.
%     \item[] Guidelines:
%     \begin{itemize}
%         \item The answer NA means that there is no societal impact of the work performed.
%         \item If the authors answer NA or No, they should explain why their work has no societal impact or why the paper does not address societal impact.
%         \item Examples of negative societal impacts include potential malicious or unintended uses (e.g., disinformation, generating fake profiles, surveillance), fairness considerations (e.g., deployment of technologies that could make decisions that unfairly impact specific groups), privacy considerations, and security considerations.
%         \item The conference expects that many papers will be foundational research and not tied to particular applications, let alone deployments. However, if there is a direct path to any negative applications, the authors should point it out. For example, it is legitimate to point out that an improvement in the quality of generative models could be used to generate deepfakes for disinformation. On the other hand, it is not needed to point out that a generic algorithm for optimizing neural networks could enable people to train models that generate Deepfakes faster.
%         \item The authors should consider possible harms that could arise when the technology is being used as intended and functioning correctly, harms that could arise when the technology is being used as intended but gives incorrect results, and harms following from (intentional or unintentional) misuse of the technology.
%         \item If there are negative societal impacts, the authors could also discuss possible mitigation strategies (e.g., gated release of models, providing defenses in addition to attacks, mechanisms for monitoring misuse, mechanisms to monitor how a system learns from feedback over time, improving the efficiency and accessibility of ML).
%     \end{itemize}
    
% \item {\bf Safeguards}
%     \item[] Question: Does the paper describe safeguards that have been put in place for responsible release of data or models that have a high risk for misuse (e.g., pretrained language models, image generators, or scraped datasets)?
%     \item[] Answer: \answerYes{} % Replace by \answerYes{}, \answerNo{}, or \answerNA{}.
%     \item[] Justification: Appendix~\ref{app:data} details the data collection process, including the filtering of risky content.
%     \item[] Guidelines:
%     \begin{itemize}
%         \item The answer NA means that the paper poses no such risks.
%         \item Released models that have a high risk for misuse or dual-use should be released with necessary safeguards to allow for controlled use of the model, for example by requiring that users adhere to usage guidelines or restrictions to access the model or implementing safety filters. 
%         \item Datasets that have been scraped from the Internet could pose safety risks. The authors should describe how they avoided releasing unsafe images.
%         \item We recognize that providing effective safeguards is challenging, and many papers do not require this, but we encourage authors to take this into account and make a best faith effort.
%     \end{itemize}

% \item {\bf Licenses for existing assets}
%     \item[] Question: Are the creators or original owners of assets (e.g., code, data, models), used in the paper, properly credited and are the license and terms of use explicitly mentioned and properly respected?
%     \item[] Answer: \answerYes{} % Replace by \answerYes{}, \answerNo{}, or \answerNA{}.
%     \item[] Justification: We have cited the original owners of assets used in the paper properly.
%     \item[] Guidelines:
%     \begin{itemize}
%         \item The answer NA means that the paper does not use existing assets.
%         \item The authors should cite the original paper that produced the code package or dataset.
%         \item The authors should state which version of the asset is used and, if possible, include a URL.
%         \item The name of the license (e.g., CC-BY 4.0) should be included for each asset.
%         \item For scraped data from a particular source (e.g., website), the copyright and terms of service of that source should be provided.
%         \item If assets are released, the license, copyright information, and terms of use in the package should be provided. For popular datasets, \url{paperswithcode.com/datasets} has curated licenses for some datasets. Their licensing guide can help determine the license of a dataset.
%         \item For existing datasets that are re-packaged, both the original license and the license of the derived asset (if it has changed) should be provided.
%         \item If this information is not available online, the authors are encouraged to reach out to the asset's creators.
%     \end{itemize}

% \item {\bf New assets}
%     \item[] Question: Are new assets introduced in the paper well documented and is the documentation provided alongside the assets?
%     \item[] Answer: \answerNA{} % Replace by \answerYes{}, \answerNo{}, or \answerNA{}.
%     \item[] Justification: This paper does not release new assets. While we provide a detailed account of our dataset creation process in Appendix~\ref{app:data}, the dataset is not publicly released at this stage due to data ownership restrictions. In addition, the model weights will be made available after the double-blind review process.
%     \item[] Guidelines:
%     \begin{itemize}
%         \item The answer NA means that the paper does not release new assets.
%         \item Researchers should communicate the details of the dataset/code/model as part of their submissions via structured templates. This includes details about training, license, limitations, etc. 
%         \item The paper should discuss whether and how consent was obtained from people whose asset is used.
%         \item At submission time, remember to anonymize your assets (if applicable). You can either create an anonymized URL or include an anonymized zip file.
%     \end{itemize}

% \item {\bf Crowdsourcing and research with human subjects}
%     \item[] Question: For crowdsourcing experiments and research with human subjects, does the paper include the full text of instructions given to participants and screenshots, if applicable, as well as details about compensation (if any)? 
%     \item[] Answer: \answerNA{} % Replace by \answerYes{}, \answerNo{}, or \answerNA{}.
%     \item[] Justification: This paper does not involve crowdsourcing nor research with human subjects.
%     \item[] Guidelines:
%     \begin{itemize}
%         \item The answer NA means that the paper does not involve crowdsourcing nor research with human subjects.
%         \item Including this information in the supplemental material is fine, but if the main contribution of the paper involves human subjects, then as much detail as possible should be included in the main paper. 
%         \item According to the NeurIPS Code of Ethics, workers involved in data collection, curation, or other labor should be paid at least the minimum wage in the country of the data collector. 
%     \end{itemize}

% \item {\bf Institutional review board (IRB) approvals or equivalent for research with human subjects}
%     \item[] Question: Does the paper describe potential risks incurred by study participants, whether such risks were disclosed to the subjects, and whether Institutional Review Board (IRB) approvals (or an equivalent approval/review based on the requirements of your country or institution) were obtained?
%     \item[] Answer: \answerNA{} % Replace by \answerYes{}, \answerNo{}, or \answerNA{}.
%     \item[] Justification: This paper does not involve crowdsourcing nor research with human subjects.
%     \item[] Guidelines:
%     \begin{itemize}
%         \item The answer NA means that the paper does not involve crowdsourcing nor research with human subjects.
%         \item Depending on the country in which research is conducted, IRB approval (or equivalent) may be required for any human subjects research. If you obtained IRB approval, you should clearly state this in the paper. 
%         \item We recognize that the procedures for this may vary significantly between institutions and locations, and we expect authors to adhere to the NeurIPS Code of Ethics and the guidelines for their institution. 
%         \item For initial submissions, do not include any information that would break anonymity (if applicable), such as the institution conducting the review.
%     \end{itemize}

% \item {\bf Declaration of LLM usage}
%     \item[] Question: Does the paper describe the usage of LLMs if it is an important, original, or non-standard component of the core methods in this research? Note that if the LLM is used only for writing, editing, or formatting purposes and does not impact the core methodology, scientific rigorousness, or originality of the research, declaration is not required.
%     %this research? 
%     \item[] Answer: \answerYes{} % Replace by \answerYes{}, \answerNo{}, or \answerNA{}.
%     \item[] Justification: We evaluate our proposed large language model on standard benchmarks and present example dialogues with the model in Sec.~\ref{sec:case} and Appendix~\ref{app:more_case}.
%     \item[] Guidelines:
%     \begin{itemize}
%         \item The answer NA means that the core method development in this research does not involve LLMs as any important, original, or non-standard components.
%         \item Please refer to our LLM policy (\url{https://neurips.cc/Conferences/2025/LLM}) for what should or should not be described.
%     \end{itemize}

% \end{enumerate}







%%%%%%%%%%%%%%%%%%%%%%%%%%%%%%%%%%%%%%%%%%%%%%%%%%%%%%%%%%%%
\newpage
\appendix

\section{Formulation of Masked Diffusion Models}
\label{app:formulation}






\begin{algorithm}[t!]
    \caption{Pre-training of LLaDA}
    \label{alg:pretrain}
    \begin{algorithmic}[1]
        \REQUIRE mask predictor $p_{\theta}$, data distribution $\pdata$
        \REPEAT
        \STATE $x_0 \sim \pdata$ \hfill \# with a probability of 1\%, the sequence length of \( x_0 \) follows \( \text{U}[1, 4096] \)
        \STATE $t \sim \text{U}(0, 1]$
        \STATE $x_t \sim q_{t|0}(x_t|x_0)$ \hfill \# $q_{t|0}$ is defined in Eq.~(\ref{eq:forward-app})
        \STATE Calculate $ \mathcal{L} = -  \frac{1}{t * L} \sum_{ i = 1 }^L \textbf{1}[x_t^i = \textrm{M}] \log p_{\theta}(x_0^i|x_t) $ \hfill \# $L$ is the sequence length of $x_0$
        \STATE Calculate $\nabla_\theta \mathcal{L}$ and run optimizer.
        \UNTIL Converged
        \STATE \textbf{Return} $p_{\theta}$
    \end{algorithmic}
\end{algorithm}


\begin{algorithm}[t!]
    \caption{Supervised Fine-Tuning of LLaDA}
    \label{alg:sft}
    \begin{algorithmic}[1]
        \REQUIRE mask predictor $p_{\theta}$, pair data distribution $\pdata$
        \REPEAT
        \STATE $p_0, r_0 \sim \pdata$ \hfill \# please refer to Appendix~\ref{app:data} for details about the SFT dat
        \STATE $t \sim \text{U}(0, 1]$
        \STATE $r_t \sim q_{t|0}(r_t|r_0)$ \hfill \# $q_{t|0}$ is defined in Eq.~(\ref{eq:forward-app})
        \STATE Calculate $ \mathcal{L} = -  \frac{1}{t * L'} \sum_{ i = 1 }^{L'} \textbf{1}[r_t^i = \textrm{M}] \log p_{\theta}(r_0^i|p_0, r_t) $ \hfill \# $L'$ is the sequence length of $r_0$
        \STATE Calculate $\nabla_\theta \mathcal{L}$ and run optimizer.
        \UNTIL Converged
        \STATE \textbf{Return} $p_{\theta}$
    \end{algorithmic}
\end{algorithm}

\begin{algorithm}[t!]
    \caption{Conditional Log-likelihood Evaluation of LLaDA}
    \label{alg:likelihood}
    \begin{algorithmic}[1]
        \REQUIRE mask predictor $p_{\theta}$, prompt $p_0$, response $r_0$, the number of Monte Carlo estimations $n_{mc}$
        \STATE $\text{log}\_\text{likelihood}=0 $
            \FOR{$i \gets 1$ to $n_{mc}$}
                \STATE $l \sim \{1, 2, \dots, L\}$ \hfill \# $L$ is the sequence length of $r_0$
                \STATE Obtain \( r_l \) by uniformly sampling \( l \) tokens from \( r_0 \) without replacement for masking
                \STATE $\text{log}\_\text{likelihood} = \text{log}\_\text{likelihood} + \frac{L}{l} \sum_{i=1}^L \textbf{1}[r_l^i = \textrm{M}] \log p_{\theta}(r_0^i|p_0, r_l) $
            \ENDFOR
        \STATE $\text{log}\_\text{likelihood} = \text{log}\_\text{likelihood} / n_{mc} $
        \STATE \textbf{Return} $\text{log}\_\text{likelihood}$
    \end{algorithmic}
\end{algorithm}
 

\subsection{Training}
\label{app:for-train}

MDMs~\citep{austin2021structured,lou2023discrete,shi2024simplified,sahoo2024simple,ou2024your} define the model distribution \( p_\theta(x_0) \) in a manner distinct from autoregressive models.

These models introduce a forward process \( \{x_t\} \) indexed by a time \( t \in [0, 1] \). This process gradually and independently masks all tokens in the sequence \( x_0 \). At time \( t = 0 \), the data point \( x_0 \) is fully observed with no masks, while for \( t \in (0, 1] \), \( x_t \) represents latent variables with varying mask ratios in expectation.

Formally, the conditional distribution of \( x_t \) given \( x_0 \) is defined by a fully factorized form:
\begin{align}
\label{eq:forward-app}
    q_{t|0}(x_t|x_0) = \prod_{i=1}^{L} q_{t|0}(x_t^i|x_0^i),
\end{align}
where the conditional distribution for each token is given by:
\begin{align}
    q_{t|0}(x_t^i|x_0^i) = 
    \begin{cases}
        1 - t, & x_t^i = x_0^i, \\
        t, & x_t^i = \textrm{M}.
    \end{cases}
\end{align}



\begin{algorithm}[t!]
    \caption{Random Remasking Strategy of LLaDA}
    \label{alg:reverse}
    \begin{algorithmic}[1]
        \REQUIRE mask predictor $p_{\theta}$, prompt $p_0$, answer length $L$, sampling steps $N$
        \STATE Set \( r_1 \) is a fully masked sequence of length \( L \).
        \FOR{$t \gets 1$ \textbf{down to} $\frac{1}{N}$ \textbf{step} $\frac{1}{N}$}
            \STATE $s = t - \frac{1}{N}$
            \STATE $r_0 = \arg\max_{r_0} p_{\theta}(r_0|p_0, r_t)$ \hfill \# we employ greedy sampling when predicting masked tokens
            \FOR{$i \gets 1$ to $L$}
                \IF{$r_t^i \neq \textrm{M}$}
                    \STATE $r_0^i = r_t^i$
                \ELSE
                    \STATE with probability \( \frac{s}{t} \), \( r_0^i \) is set to \( \textrm{M} \)
                \ENDIF
            \ENDFOR
        \STATE $r_s = r_0$
        \ENDFOR
        \STATE \textbf{Return} $r_0$
    \end{algorithmic}
\end{algorithm}


Here, \( \textrm{M} \) denotes the mask token. Intuitively, each token either remains unchanged or is masked, with the probability of being masked increasing linearly as \( t \) progresses from \( 0 \) to \( 1 \). At \( t = 1 \), all tokens are guaranteed to be masked, meaning that \( x_1 \) follows a Dirac distribution concentrated on a sequence of fully masked tokens. Notably, the linear masking probability is analogous to but distinct from, the noise schedule in continuous diffusion models~\cite{sohl2015deep,ho2020denoising,song2020score}. This linearity is motivated by the assumption that the information in the text is proportional to the number of tokens on average, making it reasonable to lose information linearly during the forward process. 

The forward process is not only reversible but also corresponds to a reverse process that is fully factorized across all tokens. The reverse process, from time \( t = 1 \) to \( 0 \), generates new data from sequences of fully masked tokens. The conditional distribution for the reverse process, for \( 0 \leq s < t \leq 1 \), is factorized as:
\begin{align}
\label{eq:reverse_process}
    q_{s|t}(x_s|x_t) = \prod_{i=1}^{L} q_{s|t}(x_s^i|x_t),
\end{align}
where the conditional distribution for each token is:
\begin{align}
\label{eq:x0_prediction}
    q_{s|t}(x_s^i|x_t) = 
    \begin{cases}
        1, & x_t^i \neq \textrm{M}, \, x_s^i = x_t^i, \\
        \frac{s}{t}, & x_t^i = \textrm{M}, \, x_s^i = \textrm{M}, \\
        \frac{t - s}{t} q_{0|t}(x_s^i|x_t), & x_t^i = \textrm{M}, \, x_s^i \neq \textrm{M}, \\
        0, & \textrm{otherwise}.
    \end{cases}
\end{align}
Thus, the key function to estimate is the conditional distribution \( q_{0|t}(x_s^i|x_t) \), 
which predicts the original token if it is masked in the input \( x_t \). This is analogous to the \emph{data prediction} form in continuous diffusion models. 
 
As proven in~\cite{ou2024your}, an equivalent yet \emph{time-free} parameterization can be derived as:
\begin{align} 
    q_{0|t}(x_s^i|x_t) = p_{\textrm{data}}(x_0^{i}|x_t^{\textrm{UM}}), \quad \forall i \textrm{ such that } x_t^{i} = \textrm{M},
\end{align}
where \( x_t^{\textrm{UM}} \) denotes the collection of unmasked tokens in \( x_t \), 
which is identical to the corresponding tokens in the original data \( x_0 \) since unmasked tokens are solely determined by \( x_0 \) and are independent of time \( t \). Intuitively, this implies that estimating the data prediction function is equivalent to estimating the conditional distributions on clean data, which is time-invariant. Consequently, the time $t$ need not be provided as input to the parametric model.

Although the development of masked diffusion is nontrivial, the implementation is straightforward. We first introduce the \emph{mask predictor}, a parametric model \( p_{\theta}(\cdot|x_t) \) (e.g., a Transformer without causal mask), which takes \( x_t \) for any $t$ as input and predict all masked tokens simultaneously. Then, we define the model distribution $p_\theta(x_0)$ as follows: starting with \( x_1 \) as a sequence of fully masked tokens, we simulate an approximate reverse process parameterized by \(  p_{\theta}(\cdot|x_t) \) from \( t = 1 \) to \( 0 \). The marginal distribution induced at \( t = 0 \) then represents the model distribution $p_\theta(x_0)$.

Formally, the mask predictor is trained using a cross-entropy loss with masking:
\begin{align}
\label{eq:objective-app}
   \mathcal{L}(\theta) \triangleq -  \mathbb{E}_{t, x_0,  x_t}  \left[\frac{1}{t} \sum_{ i = 1 }^L \textbf{1}[x_t^i = \textrm{M}] \log p_{\theta}(x_0^i|x_t) \right], 
\end{align} 
where \( x_0 \) is sampled from the training data, \( t \) is sampled uniformly from \( [0, 1] \), and \( x_t \) is sampled from \( q_{t|0}(x_t| x_0) \). The indicator function \( \textbf{1}[\cdot] \) ensures that the cross-entropy loss is computed only for masked tokens. In~\citet{shi2024simplified, sahoo2024simple, ou2024your}, it has been proven that the loss function \( \mathcal{L}(\theta) \) is an upper bound on the negative log-likelihood of the model distribution:
\begin{align}
    - \mathbb{E}_{x_0\sim p_{\textrm{data}(x_0)}} \left[\log p_\theta(x_0) \right ]\le  \mathcal{L}(\theta).
\end{align}
In summary, this principled approach trains a generative model by progressively masking tokens during a forward process and learning to recover the data distribution during a reverse process, all under the (approximate) maximum likelihood estimation framework.

\begin{algorithm}[t!]
    \caption{Low-confidence Remasking Strategy of LLaDA}
    \label{alg:low-confidence-remask}
    \begin{algorithmic}[1]
        \REQUIRE mask predictor $p_{\theta}$, prompt $p_0$, answer length $L$, sampling steps $N$
        \STATE Set \( r_1 \) is a fully masked sequence of length \( L \).
        \FOR{$t \gets 1$ \textbf{down to} $\frac{1}{N}$ \textbf{step} $\frac{1}{N}$}
            \STATE $s = t - \frac{1}{N}$
            \FOR{$i \gets 1$ to $L$}
                \IF{$r_t^i \neq \textrm{M}$}
                    \STATE $r_0^i=r_t^i$, $c^i=1$
                \ELSE
                    \STATE $r_0^i = \arg\max_{r_0^i} p_{\theta}(r_0^i|p_0, r_t)$
                    \STATE $c^i = p_{\theta}(r_0^i|p_0, r_t)_{r_0^i}$
                \ENDIF
            \ENDFOR
            \STATE  $n_{un}=\lfloor L (1 - s) \rfloor$ \hfill \# the number of unmasked tokens is $n_{un}$ in timestep $s$
            \FOR{$i \gets 1$ to $L$}
                \IF{$c^i \in \text{Lowest}-n_{un} \left(\{c^i \}_1^{L} \right)$}
                    \STATE $r_0^i = \textrm{M}$ \hfill \# the $n_{un}$ positions with the least confidence are selected for remasking.
                \ENDIF
            \ENDFOR
        \STATE $r_s = r_0$
        \ENDFOR
        \STATE \textbf{Return} $r_0$
    \end{algorithmic}
\end{algorithm}


\subsection{Inference}
\label{app:for-infer}
   
The cross-entropy loss in Eq.~(\ref{eq:objective-app}) has several equivalent forms~\cite{ou2024your}. The first one is given by
\begin{align}
    -\mathbb{E}_{l\sim \{1, 2, \dots, L\}, x_0, x_l} \left[\frac{L}{l}\sum_{i=1}^L \textbf{1}[x_l^i = \textrm{M}] \log p_{\theta}(x_0^i|x_l) \right],    
\label{eq:ppl-eval-app}
\end{align}
where \( l \) is uniformly sampled from \( \{1, 2, \dots, L\} \), 
and \( x_l \) is obtained by uniformly sampling \( l \) tokens from \( x_0 \) without replacement for masking. Despite masking exactly $l$ tokens is different from masking each token independently with probability $t$, these two masking methods lead to equivalent results in expectation \cite{ou2024your}.


While Eq.~(\ref{eq:objective-app}) and Eq.~(\ref{eq:ppl-eval-app}) share the same expectation, their variances differ. Intuitively, in Eq.~(\ref{eq:objective-app}), we expect \(x_t\) to have a fraction of \(t\) tokens masked. However, the randomness of the forward process (i.e., Eq.~(\ref{eq:forward-app})) often causes deviations, especially when \(x_t\) contains few tokens. In contrast, in Eq.~(\ref{eq:ppl-eval-app}), the fraction of masked tokens in \(x_l\) is deterministically \(\frac{l}{L}\). While a theoretical analysis depends on the data distribution, empirical results show that Eq.~(\ref{eq:objective-app}) requires over 1000 Monte Carlo estimates for stable results, whereas Eq.~(\ref{eq:ppl-eval-app}) achieves stability with only 128 estimates. In addition, we can simply modify Eq.~(\ref{eq:ppl-eval-app}) to its conditional version (i.e., Eq.~(\ref{eq:ppl-eval})) based on Eq.~(\ref{eq:sft-objective}).


Any-order autoregressive models (AO-ARM)~\cite{hoogeboom2021autoregressive, UriaML14,Shih2022TrainingAI} characterize the joint distribution autoregressively for all possible orders $\pi$ of the $L$ variables. To learn such a distribution, an AO-ARM utilizes a weight-sharing neural network to model all univariate conditionals and employs mask tokens to represent absent variables. During training, the expected negative log-likelihood over the uniform distribution of all orders $U_\pi$ is minimized:
\begin{align}
\label{eq:ao_obj}
-\mathbb{E}_{x_0, \pi \sim U_\pi}\left[ \sum_{i=1}^L \log p_\theta(x_0^{\pi(i)}|x_0^{\pi(<i)};\pi)\right].
\end{align}
Intuitively, $x_0^{\pi(<i)}$ can be understood as a masked token $x_t$ with index in ${\pi(\geq i)}$ being masked. It can be further proved that 
Eq.~(\ref{eq:ao_obj}) is equivalent to Eq.~(\ref{eq:objective-app}). This connection explains the bidirectional reasoning capabilities of LLaDA, even though it was never used explicitly in the inference procedure.

In addition, \citet{nie2024scaling} introduces unsupervised classifier-free guidance (CFG), a plug-and-play technique that balances alignment with prompts and text diversity. Specifically, unsupervised CFG employs the following modified mask predictor for inference:
\begin{align}
    \tilde{p}_{\theta}(r_0| p_0, r_t ) \propto \frac{p_{\theta}(r_0| p_0, r_t)^{1+w}}{p_{\theta}(r_0|m, r_t)^w},
\end{align}
where $m$ is a mask sequence of the same length as $p_0$ and $w$ is a tunable hyperparameter that controls the strength of $p_0$. To ensure a fair comparison with ARMs, we do not apply CFG to LLaDA in the main text. However, as demonstrated in Appendix~\ref{app:cfg}, LLaDA is fully compatible with CFG and consistently exhibits improved performance when it is applied.

\subsection{Algorithms}
\label{app:algorithms}
In this section, we present the training and inference algorithms. Specifically, we introduce the pre-training and supervised fine-tuning algorithms in Algorithm~\ref{alg:pretrain} and Algorithm~\ref{alg:sft}, respectively. In addition, the likelihood evaluation algorithm is provided in Algorithm~\ref{alg:likelihood}. Finally, we present the reverse generation process in Algorithm~\ref{alg:reverse} and Algorithm~\ref{alg:low-confidence-remask}, which correspond to the random remasking and the low-confidence~\citep{chang2022maskgit} remasking strategy, respectively.

\section{Experiments}
\label{app:experiment}

\subsection{Data Collection and Preprocessing}
\label{app:data}
In this section, we first introduce the data collection and filtering processes for both pre-training and SFT. We then describe how LLaDA leverages these datasets during training.

Our pre-training corpus is constructed from diverse publicly available sources, including web data, books, academic papers, social media, encyclopedias, mathematics, and code, with approximately 11\% Chinese, 61\% English, and 28\% code. The data cleaning process involves PDF text extraction, deduplication, and harmful content filtering. To further ensure quality, we fine-tune a BERT~\citep{devlin2018bert} model for automated data quality annotation, enabling the selection of higher-quality samples. Our SFT dataset consists of 1 million human-annotated samples and 3.5 million synthetic samples, generated using methods similar to those proposed in~\citet{xu2024magpie, wei2023magicoder}.

We concatenate the collected documents in the pre-training corpus and segment the text into fixed-length sequences according to the predefined sequence length.

For SFT, a dynamic sequence length strategy is employed, where $|\text{EOS}|$ tokens are appended to the end of shorter pairs to ensure uniform sequence lengths across all samples within each mini-batch. Notably, the padding $|\text{EOS}|$ tokens are treated as part of the response, i.e., masked and included in the training objective. The $|\text{EOS}|$ tokens are removed from the generated outputs during sampling. This strategy ensures that the model learns to control the length of its responses by generating $|\text{EOS}|$, enabling the response length to align effectively with the given prompt.

In addition, for $n$-turn dialogues $(p_0^0, r_0^0, p_0^1, r_0^1, \dots, p_0^{n-1}, r_0^{n-1})$, we treat it as $n$ single-turn dialogue pairs, i.e., $(p_0^0, r_0^0), (p_0^0r_0^0p_0^1, r_0^1), \dots, (p_0^0r_0^0p_0^1r_0^1\dots p_0^{n-1}, r_0^{n-1})$ and randomly sample one. This data partitioning strategy not only equips LLaDA with multi-turn dialogue capabilities but also aligns with the above $|\text{EOS}|$ padding strategy.


\subsection{Details about Model Training}
\label{app:exp-1b-config}
This section provides the training details of LLaDA and the corresponding ARM baselines.

Firstly, for efficiency, we trained an ARM and an MDM, both with 1.5B parameters and identical architectures. Additionally, we scaled the MDM to 8B parameters. Due to computational resource constraints, we did not train an 8B autoregressive model with the same architecture. Instead, we utilized our previously trained 7B autoregressive model for comparison. These four models are utilized in the scalability analysis in Sec.~\ref{sec:scaling}.

We adopted a Transformer architecture similar to LLaMA~\cite{dubey2024llama, touvron2023llama2} for the ARMs and MDMs we trained. Specifically, we employ RMSNorm~\cite{zhang2019root} to stabilize training, use SwiGLU~\cite{shazeer2020glu} as the activation function to enhance non-linearity, and integrate RoPE~\cite{su2024roformer} for more expressive positional encoding. Tab.~\ref{table:8b} provides an overview of the model architectures.

For the 1B and 7B ARM baselines, as well as the 1B and 8B LLaDA models, we utilized the AdamW optimizer~\citep{loshchilov2017decoupled} with a weight decay of 0.1 and adopted the Warmup-Stable-Decay~\citep{hu2024minicpm} learning rate scheduler. The learning rate was linearly increased from 0 to the maximum value over the first 2000 iterations and then held constant. For LLaDA 8B, to ensure stable training, the learning rate was reduced once during pre-training, as detailed in Sec.~\ref{sec:pre-traing}. For the 1B ARM baseline and both the 1B and 8B LLaDA models, the maximum learning rate is set to $4 \times 10^{-4}$ with a batch size of 1280, without any hyperparameter tuning. For the 7B ARM baseline, the maximum learning rate is set to $4.2 \times 10^{-4}$ with a batch size of 4224, both selected via grid search.

Additionally, we employ the widely used $6ND$ formulation~\cite{kaplan2020scaling, hoffmann2022training} to calculate the training FLOPs in Fig.~\ref{fig:scaling}, where $N$ represents the number of non-embedding parameters, and $D$ denotes the total number of training tokens. The detailed results corresponding to Fig.~\ref{fig:scaling} are provided in Tab.~\ref{tab:scaling_llada} and Tab.~\ref{tab:scaling_arm}.


 \begin{table}[t!]
    \centering
    \caption{\textbf{Model Architecture.} We report the architectural configurations for our 1B and 7B ARM baselines, the 1B and 8B LLaDA models, and the 8B LLaMA3 model. }
    \vspace{.2cm}
    \label{table:8b}
    \begin{adjustbox}{max width=\textwidth}
    \begin{tabular}{lccccc}
      \toprule
         & Our ARM Baseline 1B & LLaDA 1B & Our ARM Baseline 7B & LLaDA 8B & LLaMA3 8B \\
         \midrule
         Layers & 22 & 22 & 28 & 32 & 32 \\
         \midrule
         Model dimension & 2048  & 2048 & 4096 &4096 & 4096 \\
         \midrule
         Attention heads & 32 & 32 & 32 & 32 & 32 \\
         \midrule
         Vocabulary size & 126,464 & 126,464 & 126,464 & 126,464 & 128,000 \\
         \midrule
         FFN dimension &5634 & 5634 & 13,440 & 12,288 & 14,336 \\   
         \midrule
         Key/Value heads & 4 & 4 & 8 & 32 & 8 \\
         \midrule
         Total parameters & 1.49 B & 1.49 B & 6.83 B & 8.02 B & 8.03 B \\
         \midrule
         Non-embedding parameters & 0.97 B & 0.97 B & 5.80 B & 6.98 B & 6.98 B\\
      \bottomrule
    \end{tabular}
    \end{adjustbox}
    \vspace{-.2cm}
\end{table}

\subsection{Ablation on Classifier-free Guidance}
\label{app:cfg}
This section presents an ablation study on classifier-free guidance (CFG). Theoretical details about CFG can be found in the Appendix~\ref{app:for-infer}.

For simplicity, we select six representative benchmarks, including ARC-C, HellaSwag, TruthfulQA, WinoGrande, PIQA, and GPQA, and conduct experiments using LLaDA 8B Base. We search the CFG scale in $\{0.5, 1, 1.5, 2\}$ for each task and report the best result. As shown in Tab.~\ref{tab:ablation-cfg}, CFG consistently improves the performance of LLaDA. We emphasize that, to ensure a fair comparison with ARMs, CFG is not used in the main results reported in the paper.

\begin{table}[t!]
    \centering
    \caption{\textbf{Ablation on CFG.} CFG consistently improves the performance of LLaDA.}
    \label{tab:ablation-cfg}
    \vspace{.2cm}
    \begin{adjustbox}{max width=\textwidth}
    \begin{tabular}{lcccccc}
      \toprule
      & ARC-C & Hellaswag & TruthfulQA & WinoGrande & GPQA & PIQA \\
      \midrule
      w/o CFG & 45.9 & 70.5 & 46.1 & \textbf{74.8} & 25.2 & 73.6 \\
      w/ CFG & \textbf{47.9} & \textbf{72.5} & \textbf{46.4} & \textbf{74.8} & \textbf{26.1} & \textbf{74.4} \\
      \bottomrule
    \end{tabular}
    \end{adjustbox}
    \vspace{-.2cm}
\end{table}





  

\subsection{Details and Ablation on Sampling Strategies}
\label{app:sample}
In this section, we first introduce the different sampling strategies supported by LLaDA. We then present ablation studies to evaluate the performance of these sampling strategies.

\begin{figure}[t]
  \centering
  \begin{subfigure}{0.25\textwidth}
    \centering
    \includegraphics[width=\linewidth]{imgs/sample/ar_sample.jpg}
    \caption{Autoregressive.}
  \end{subfigure}\hfill
  \begin{subfigure}{0.36\textwidth}
    \centering
    \includegraphics[width=\linewidth]{imgs/sample/block_diffusion.jpg}
    \caption{Block Diffusion.}
  \end{subfigure}\hfill
  \begin{subfigure}{0.36\textwidth}
    \centering
    \includegraphics[width=\linewidth]{imgs/sample/block_diffusion_llada.jpg}
    \caption{Block Diffusion LLaDA.}
  \end{subfigure}
  \caption{\textbf{Flexible Sampling Strategies Supported by LLaDA.} Colored squares depict non‑masked tokens, while squares marked with $\times$ denote masked tokens. In this illustration, the block length for both block diffusion and block diffusion LLaDA sampling is 4.}
  \label{fig:flexible_sampling}
  \vspace{-0.3cm}
\end{figure}



\begin{table}[t!]
    \centering 
    \caption{\textbf{Ablation on Sampling Strategies for LLaDA 8B Base.} $L'$ is the block length. Pure diffusion sampling achieves the best overall performance.} 
    \label{tab:abation_sample_base}
    \vspace{0.2cm}
    \begin{adjustbox}{max width=\textwidth}
    \begin{tabular}{cc|ccccc}
    \toprule
    & & BBH & GSM8K & Math & HumanEval & MBPP \\
    \midrule
    Autoregressive & & 38.1 & 63.1 & 23.6 & 18.3 & 33.4 \\
    \midrule
    \multirow{4}{*}{Block Difusion} &$L' = 2$ &37.3 & 62.6 & 25.2 & 14.6 & 33.6 \\
    & $L'= 4$ &40.0 & 65.7 & 26.6 & 15.9 & 36.0 \\
    & $L' = 8$ &42.0 & 68.2 & 27.7 & 19.5 & 39.2 \\  
    & $L' = 32$ & 45.7 & 68.6 & 29.7 & 29.9 & 37.4 \\
    \midrule
    \multirow{4}{*}{Block Diffusion LLaDA} & $L' = 2$ & 48.0 & 70.0 & 30.8 & 26.2 & \textbf{40.0} \\
    & $L' = 4$ & 48.5 & \textbf{70.3} & 31.3 & 27.4 & 38.8 \\
    & $L' = 8$ & 48.6 & 70.2 & 30.9 & 31.1 & 39.0 \\
    & $L' = 32$ & 48.3 & \textbf{70.3} & 31.2 & 32.3 & \textbf{40.0} \\
    \midrule
    Pure Diffusion & & \textbf{49.7} & \textbf{70.3} & \textbf{31.4} & \textbf{35.4} & \textbf{40.0}\\
    \bottomrule
    \end{tabular}
    \end{adjustbox}
\end{table}


\begin{table}[t!]
    \centering 
    \caption{\textbf{Ablation on Sampling Strategies for LLaDA 8B Instruct.} The block length is set to 32 for efficiency. Pure diffusion sampling achieves the best overall performance.} 
    \label{tab:abation_sample_instruct}
    \vspace{0.2cm}
    \begin{adjustbox}{max width=\textwidth}
    \begin{tabular}{c|ccccccc}
    \toprule
    & GSM8K & Math & HumanEval & MBPP & GPQA & MMLU-Pro & ARC-C \\
    \midrule
    Autoregressive & 0 & 9.5 & 0 &0 & 0 & 0 & 84.4\\
    \midrule
    Block Diffusion & 24.6 & 23.5 & 17.1 & 21.2& 29.3 & 32.5 & 88.1\\
    \midrule
    Block Difusion LLaDA & \textbf{77.5} & \textbf{42.2} & 46.3 & 34.2 & 31.3 & 34.8 & 85.4 \\
    \midrule
    Pure Diffusion & 69.4 & 31.9 & \textbf{49.4} & \textbf{41.0} & \textbf{33.3}  & \textbf{37.0} & \textbf{88.5}\\
    \bottomrule
    \end{tabular}
    \end{adjustbox}
\end{table}

\textbf{Flexible Sampling Strategies.} In Sec.~\ref{sec:inference}, Fig.~\ref{fig:overview}~(c) illustrates the reverse generation process of LLaDA. As shown in Fig.~\ref{fig:flexible_sampling}, in addition to the reverse generation process, LLaDA also supports autoregressive and block diffusion~\citep{arriola2025block} sampling directly after the pre-training or SFT stages, without requiring any further modifications or retraining. Block diffusion sampling applies the origin reverse process within each block and the autoregressive sampling across blocks. In the original block diffusion process, the sequence length varies dynamically. As shown in Fig.~\ref{fig:flexible_sampling}~(c), LLaDA can also adopt a fixed-length block diffusion strategy, which we refer to as block diffusion LLaDA, also known as semi-autoregressive remasking.


\textbf{Experimental Setup.} We evaluate different sampling strategies using both LLaDA 8B Base and LLaDA 8B Instruct for comprehensive analysis. For LLaDA 8B Base, we use the five benchmarks in Tab.~\ref{tab:base} that are evaluated based on sampling rather than likelihood estimation. For LLaDA 8B Instruct, we use the seven metrics in Tab.~\ref{tab:chat}, excluding MMLU and HellaSwag, since these two tasks only require the model to generate a single token (i.e., A, B, C, or D). In all settings, one token is generated per sampling step. For autoregressive and block diffusion sampling, generation terminates when the $|\text{EOS}|$ token is produced. For block diffusion LLaDA (i.e., semi-autoregressive remasking) and pure diffusion sampling, the generation length is fixed at 1024 for LLaDA 8B Base, while for LLaDA 8B Instruct, it is tuned from \{64, 256, 512\} to balance efficiency and performance. Low-confidence remasking is applied to intra-block diffusion sampling in both block diffusion and block diffusion LLaDA, as well as to pure diffusion sampling. We also test different block lengths for LLaDA 8B Base. For LLaDA 8B Instruct, we only evaluate block length 32 for efficiency, as it yields the best results on LLaDA 8B Base.

Additionally, for LLaDA 8B Instruct, due to heavy padding of $|\text{EOS}|$ tokens in the SFT data (as detailed in Sec.~\ref{app:data}), we observe that under pure diffusion sampling, the proportion of $|\text{EOS}|$ tokens in model outputs becomes very high. This leads to extremely short generations and degrades model performance. To mitigate this issue, for HumanEval, MBPP, GSM8K, Math, and GPQA, we set the confidence score of the $|\text{EOS}|$ token to zero during pure diffusion sampling. This adjustment helps maintain an appropriate ratio of $|\text{EOS}|$ tokens during generation.

Finally, we conduct ablation studies to analyze the effects of random and low-confidence remasking strategies using the pure diffusion sampling. For efficiency, we use LLaDA 8B Base with generation length and sampling steps set to 256 in this analysis.

\textbf{Results.} As shown in Tab.~\ref{tab:abation_sample_base}, for block diffusion sampling, overall performance improves as the block length increases. Moreover, both Tab.~\ref{tab:abation_sample_base} and Tab.~\ref{tab:abation_sample_instruct} show that block diffusion sampling consistently outperforms autoregressive sampling, and block diffusion LLaDA sampling further improves upon standard block diffusion sampling. Finally, pure diffusion sampling achieves the best overall performance.

In addition, Tab.~\ref{tab:ablation-remasking} shows that the low-confidence remasking strategy consistently outperforms the random remasking strategy. We hypothesize that low-confidence remasking functions similarly to the annealed sampling method used by default in ARMs, improving accuracy by reducing the diversity of generated sentences.

\begin{table}[t!]
    \centering
    \caption{\textbf{Analysis on Random and Low-confidence Remasking Strategies.} The low-confidence remasking consistently outperforms the random remasking.}
    \label{tab:ablation-remasking}
    \vspace{.2cm}
    \begin{adjustbox}{max width=\textwidth}
    \begin{tabular}{lccccc}
      \toprule
         Length & BBH & GSM8K & Math & HumanEval & MBPP \\
         \midrule
        Random Remasking & 32.1 &21.3 & 9.2 &11.6 & 21.0\\
        Low-confidence Remasking & \textbf{45.0} & \textbf{70.0} & \textbf{30.3} & \textbf{32.9} & \textbf{40.2}\\
      \bottomrule
    \end{tabular}
    \end{adjustbox}
    \vspace{-.2cm}
\end{table}

We discover that autoregressive sampling leads to very poor performance for LLaDA 8B Instruct. This is because each SFT data is a complete sentence, so given a sequence length, LLaDA 8B Instruct tends to generate a full sentence within that length. In contrast, LLaDA 8B Base does not suffer from this issue, as the pre-training data consists of truncated documents (as detailed in Appendix~\ref{app:data}) and the model is trained with random sequence lengths (as detailed in Sec.~\ref{sec:pre-traing}). As a result, when given a short sequence length, LLaDA 8B Base tends to generate only part of a sentence, which can then be used as a prefix to continue generation.

Setting the block length to 8 in Tab.~\ref{tab:abation_sample_instruct} further improves the GSM8K score from 77.5 to 78.6.



\subsection{Ablation on Generated Length}
\label{app:ablation-length}

\begin{table}[t!]
    \centering
    \caption{\textbf{Ablation on Generation Length.} The results are not sensitive to the length hyperparameter.}
    \label{tab:ablation-length}
    \vspace{.2cm}
    \begin{adjustbox}{max width=\textwidth}
    \begin{tabular}{lccccc}
      \toprule
         Length & BBH & GSM8K & Math & HumanEval & MBPP \\
         \midrule
        256 & 45.0	& 70.0 & 30.3& 32.9& \textbf{40.2}\\
        512 & \textbf{50.4}	& \textbf{70.8}	& 30.9	& 32.9 & 39.2\\
        1024 & 49.7	& 70.3	& \textbf{31.4}& \textbf{35.4}&	40.0\\
      \bottomrule
    \end{tabular}
    \end{adjustbox}
    \vspace{-.2cm}
\end{table}

In this section, we conduct ablation studies on the generated length. 

To ensure fairness, for each setting, we set the number of sampling steps equal to the generated length, ensuring that in each sampling step, just one tokens are transferred from the mask to the text. We conduct experiments using LLaDA 8B Base.


As reported in Tab.~\ref{tab:ablation-length}, the results are not sensitive to the length hyperparameter.




\subsection{Standard Benchmarks and Evaluation Details}
\label{app:bench}

\begin{figure}[t!]
    \centering
    \begin{minipage}{0.4\textwidth}
        \centering
        \includegraphics[width=\textwidth]{imgs/efficiency/GSM8K_efficiency.pdf}
    \end{minipage}
    \hspace{0.1\textwidth}
    \begin{minipage}{0.4\textwidth}
        \centering
        \includegraphics[width=\textwidth]{imgs/efficiency/Math_efficiency.pdf}
    \end{minipage}
    \begin{minipage}{0.4\textwidth}
        \centering
        \includegraphics[width=\textwidth]{imgs/efficiency/HumanEval_efficiency.pdf}
    \end{minipage}
    \hspace{0.1\textwidth}
    \begin{minipage}{0.4\textwidth}
        \centering
        \includegraphics[width=\textwidth]{imgs/efficiency/MBPP_efficiency.pdf}
    \end{minipage}
    \caption{\textbf{Analysis of Sampling Efficiency.} The generation length for LLaDA is set to 256, with sampling steps set to 32, 64, 128, and 256 across the figures. This corresponds to decoding 8, 4, 2, and 1 token(s) per forward pass, respectively. LLaDA enables a flexible trade-off between generation quality and sampling speed. }
    \label{fig:efficiency-analysis}
    \vspace{-.2cm}
\end{figure}

In this section, we introduce the benchmarks we used and present the details of our evaluation process.

Following standard LLM~\cite{qwen2, qwen2.5} evaluation practices, we assess LLaDA across four dimensions:

\textbf{General ability:} MMLU~\cite{hendrycks2020measuring}, BBH~\cite{suzgun2022challenging}, ARC-C~\cite{clark2018think}, Hellaswag~\cite{zellers2019hellaswag}, TruthfulQA~\cite{lin2021truthfulqa}, WinoGrande~\cite{sakaguchi2021winogrande} and PIQA~\cite{bisk2020piqa}.

\textbf{Math and science ability:} GSM8K~\cite{cobbe2021training}, Math~\cite{hendrycks2021measuring} and GPQA~\cite{rein2023gpqa}.

\textbf{Code generation:} HumanEval~\cite{chen2021evaluating}, HumanEval-FIM~\cite{bavarian2022efficient} and MBPP~\cite{austin2021program}.

\textbf{Chinese understanding:} CMMLU~\cite{li2023cmmlu} and C-Eval~\cite{huang2024c}.


For all the aforementioned benchmarks, we follow the widely adopted evaluation process~\cite{eval-harness} used in LLM assessments, primarily employing conditional likelihood estimation and conditional generation. Specifically, for certain benchmarks, a prompt and multiple candidate answers are provided, and the model is required to compute each candidate's conditional likelihood. The candidate with the highest likelihood is then selected as the model’s final answer, and accuracy is used as the evaluation metric. For the remaining benchmarks, the model generates responses based on the given prompt, and performance is evaluated using metrics such as exact match and other relevant criteria.

For the base model, we use conditional likelihood estimation for MMLU, CMMLU, C-Eval, ARC-C, Hellaswag, TruthfulQA, WinoGrande, PIQA, and GPQA, while the remaining benchmarks are evaluated using conditional generation. For the instruct model, we evaluate all benchmarks using conditional generation.


For the base model, we use the widely adopted open-source evaluation framework lm-evaluation-harness~\cite{eval-harness}, except for the HumanEval-FIM metric, which is not supported by the framework. For HumanEval-FIM on the base model and all evaluation metrics on the instruct model, we use an internal evaluation library. We choose the internal library as lm-evaluation-harness shows greater deviation from the LLaMA3 results reported by~\citet{qwen2}, relative to our internal evaluation.



For benchmarks evaluated via conditional likelihood estimation, we use Monte Carlo estimation to approximate Eq.~(\ref{eq:ppl-eval}) for LLaDA. Since MMLU, CMMLU, and C-EVAL only require the likelihood of a single token, a single Monte Carlo estimate is sufficient for these benchmarks. For all other benchmarks, we find that 128 Monte Carlo samples are adequate to produce stable results. 


For benchmarks evaluated using conditional generation, we apply pure diffusion sampling with a low-confidence remasking strategy to both LLaDA Base and LLaDA Instruct. For LLaDA Base, we set both the generation length and the number of sampling steps to 1024. For LLaDA Instruct, the number of sampling steps is set equal to the answer length, which is configured as follows: 3 for MMLU and HellaSwag, 64 for GPQA, 256 for MBPP and MMLU-Pro, and 512 for HumanEval, GSM8K, Math, and ARC-C. As detailed in Appendix~\ref{app:sample}, for HumanEval, MBPP, GSM8K, Math, and GPQA, we set the confidence of the $|\text{EOS}|$ token to zero during sampling for LLaDA Instruct.





\subsection{Analysis of Sampling Efficiency}
\label{app:analysis}

In this section, we first analyze the sampling efficiency of LLaDA, including both sampling speed and memory consumption. We then discuss potential optimizations to further improve its efficiency.

We use four representative open-ended generation benchmarks for sampling speed analysis: GSM8K, Math, HumanEval, and MBPP. We use the widely adopted throughput metric to measure generation speed, defined as the number of tokens generated per second. We compare LLaDA 8B Base and LLaMA3 8B Base, both using bfloat16 precision. All experiments in this section were conducted on a single A100-80GB GPU with a batch size of 1. For LLaDA, the output length is fixed to 256 tokens across all four benchmarks.

Fig.~\ref{fig:efficiency-analysis} shows that LLaDA enables a flexible trade-off between generation quality and speed by adjusting the number of sampling steps. Specifically, on the GSM8K and Math datasets, LLaDA 8B Base achieves comparable performance to LLaMA3 8B Base while delivering 1.5 and 1.8 times higher throughput, even though LLaMA3 uses KV Cache and LLaDA operates without any inference optimization techniques.  
For the HumanEval benchmark, LLaDA 8B Base performs comparably to LLaMA3 8B Base when the throughput is matched. On the MBPP benchmark, LLaDA 8B Base lags behind LLaMA3 8B Base.

For LLaMA3, the acceleration benefit provided by KV caching is notably weaker on the HumanEval dataset, which can be attributed to its relatively short prompt lengths. Specifically, the average prompt lengths for GSM8K, Math, MBPP, and HumanEval are 894, 680, 628, and 132 tokens, respectively.

\begin{table}[t!]
    \centering 
    \caption{\textbf{Analysis of Memory Consumption.} Memory is measured in GB. Without any inference optimization techniques (e.g., KV Cache), LLaDA has memory usage comparable to LLaMA3, and slightly higher than LLaMA3 when the latter uses KV Cache.} 
    \label{tab:memory}
    \vspace{0.2cm}
    \begin{adjustbox}{max width=\textwidth}
    \begin{tabular}{cc|ccc}
    \toprule
    Input Length & Output Length & LLaDA 8B & LLaMA3 8B w/o KV-Cache & LLaMA3 8B w/ KV-Cache \\
    \midrule
    \multirow{3}{*}{512} & 512 & 17.03 & 16.70 & 16.32 \\
    & 1024 & 17.53  & 17.49 & 16.43\\
    & 2048 & 18.52  & 20.00 & 16.73 \\
    \midrule
    \multirow{3}{*}{1024} & 512 & 17.53 & 17.16 & 16.36\\
    & 1024 & 18.01 & 18.26 & 16.41\\
    & 2048 & 19.02 & 21.39 & 16.74\\
    \bottomrule     
    \end{tabular}
    \end{adjustbox}
    \vspace{-.2cm}
\end{table}


Tab.~\ref{tab:memory} compares of memory consumption between LLaDA 8B Base and LLaMA3 8B Base. To avoid variations in generation length caused by differences in training data, we fix both the input and output token lengths during the memory analysis. For LLaDA, memory usage remains constant regardless of the number of sampling steps. Its memory consumption is comparable to LLaMA3 8B Base without KV cache, but slightly higher than with KV cache.

We emphasize that the goal of this study is not to propose a model that is faster than ARMs. Instead, we aim to show the promise of diffusion models for language modeling at scale and challenge the common assumption that core LLM capabilities such as scalability, in-context learning, and instruction-following are inherently depend on ARMs. A substantial body of research~\citep{arriola2025block, xu2025show, liu2024think, zhu2025di, ren2025fast, hayakawa2024distillation, zhao2024informed, zheng2024masked, park2024jump, deschenaux2024beyond} has focused on improving the generation efficiency of MDMs through algorithmic or architectural innovations. We leave similar efficiency-oriented exploration for LLaDA to future work.









\subsection{Evaluation on iGSM Dataset}
\label{app:igsm}


\begin{table}[t!]
    \centering 
    \caption{\textbf{Comparison on iGSM Dataset.}} 
    \vspace{0.2cm}
    \begin{tabular}{lccc}
        \toprule
     & 4 steps & 5 steps & 6 steps \\ 
     \midrule
     LLaMA3 8B Base & 38.0 & 35.0 & 34.0 \\
     LLaDA 8B Base  & \textbf{64.0} & \textbf{41.0} & \textbf{44.0} \\
\bottomrule
    \end{tabular}
    \label{tbl:math}
    \vspace{-.2cm}
\end{table}

To further assess the mathematical capabilities of LLaDA, we test its performance on iGSM \cite{YXLA2024-gsm1}, 
an infinite, synthetic GSM8K-like dataset. 
iGSM is generated via specific rules, 
with parameters that control the difficulty of problems (i.e., the number of solution steps).
For evaluation consistency,
we append "\#\#\#\# \$answer" to the final solution, adhering to the GSM8K format. 
Below is an example with solution steps set to 4:


\begin{framed}
(\textbf{Question}) The number of each North Star Elementary's Cultural Studies Classroom equals 1. 
The number of each Westridge Elementary's Dance Studio equals 3 times as much as the sum of each North Star Elementary's Classroom and each North Star Elementary's Cultural Studies Classroom. 
How many Dance Studio does Westridge Elementary have? 
\\
(\textbf{Solution}) Define North Star Elementary's Cultural Studies Classroom as x; so x = 1. \\
Define North Star Elementary's Classroom as m; so m = x = 1. \\
Define Westridge Elementary's Dance Studio as n; w = m + x = 1 + 1 = 2; \\
so n = 3 * w = 3 * 2 = 1  \#\#\#\# 1
\end{framed}


Since there are slight differences between GSM8K and iGSM (e.g., the use of a mod 5 algorithmic system), we follow \cite{YXLA2024-gsm1} and provide a system prompt along with four-shot question-answer pairs for each problem.

\begin{framed}
    (\textbf{Prompt}) You're an expert at solving elementary math problems involving addition, subtraction, and multiplication. You solve all the problems in a uniform format. All calculations are done modulo 5. For example, 4 + 4 equals 3, 2 + 4 equals 1, 3 + 3 + 3 equals 4, 3 * 3 equals 4, and 2 * 2 equals 4. When providing your solution, please end with '\#\#\#\# x.' where x is your final answer, an integer between 0 and 4. You must solve all the problems using the same solution format. Our scenarios involve up to four categories of objects: schools, classrooms, backpacks and stationeries. Each school may contain classrooms, each classroom may contain backpacks, and each backpack may contain stationeries. We can specify quantities, such as \"the number of dance studios at each Lakeshore High.\" Assume that every entity with the same name has an identical configuration; for example, each Lakeshore High contains the same number of dance studios. Another guiding principle is that what is not mentioned does not exist: when we refer to classrooms at Lakeshore High, we are only discussing the classrooms explicitly mentioned in our scenario. Furthermore, if Lakeshore High is not even mentioned, any classroom within it is automatically considered to be non-existent (i.e. 0).
\end{framed}


For solution steps ranging from 4 to 6, we generate 100 questions for each case and report the corresponding accuracy in \cref{tbl:math}. As shown in the table, LLaDA 8B Base demonstrates significant and consistent advantages over LLaMA3 8B Base on unseen mathematical problems, aligning with the results in Table~\ref{tab:base}.

\subsection{Poem Completion Tasks}
\label{app:exp-poem}
In this section, we present examples from our poem completion dataset as follows.



\begin{CJK}{UTF8}{gbsn}
Example 1:\\
Prompt: 窈窕淑女的下一句是什么？直接输出句子即可。\\
Answer: 君子好逑。
\end{CJK}

\begin{CJK}{UTF8}{gbsn}
Example 2:\\
Prompt: 不拘一格降人才的上一句是什么？直接输出句子即可。\\
Answer: 我劝天公重抖擞。
\end{CJK}



\subsection{More Case Studies}
\label{app:more_case}
In this section, we present additional case studies of LLaDA 8B Instruct. First, Tab.~\ref{tab:visual-sar} shows the sampling process of the block diffusion LLaDA sampling, while Tab.~\ref{tab:visual-multi} depicts the sampling process for multi-turn dialogues with random remasking. Additionally, Tab.~\ref{tab:case-single} and Tab.~\ref{tab:case-multi} provide further examples of single-turn and multi-turn dialogues. Finally, Tab.~\ref{tab:case-reversal} presents examples of poem reversal completions where the LLaDA 8B Instruct model succeeds, in contrast to the failure of GPT-4o.


\section{Impact Statement}
\label{app:impact}
Our work shows the promise of diffusion models for language modeling at scale and challenges the common assumption that core LLM capabilities such as scalability, in-context learning, and instruction-following are inherently dependent on ARMs. Our findings open new avenues for exploring alternative probabilistic paradigms in natural language processing, with potential applications in conversational AI, code generation, and complex reasoning tasks.

However, diffusion models, like traditional LLMs, raise similar societal concerns. These include the environmental impact of large-scale training, the potential misuse for generating harmful content, and the amplification of biases present in training data. Addressing these challenges is critical to ensuring the responsible development and deployment of diffusion language models.






\begin{table}[t!]
\centering
\caption{\textbf{Visualization of the Block Diffusion LLaDA Sampling Process.} In the response of LLaDA, darker colors indicate tokens predicted in the later stages of sampling, while lighter colors correspond to earlier predictions.}
\label{tab:visual-sar}
\vspace{5pt} %
\scalebox{1}{
\begin{adjustbox}{max width=\textwidth}
\begin{tabular}{l p{12.5cm} }
\toprule
User & What are the benefits of regular exercise for physical and mental health?\\
LLaDA & \includegraphics[width=\linewidth,valign=t]{imgs/visualization/sar.jpg}\\
\bottomrule
\end{tabular}
\end{adjustbox}
}
\end{table}



\begin{table}[t!]
  \begin{minipage}{0.99\textwidth}
\centering
\caption{\textbf{Visualization of the Multi-turn Dialogue.} We employ random remasking strategy. In the response of LLaDA, darker colors indicate tokens predicted in the later stages of sampling, while lighter colors correspond to earlier predictions.}
\label{tab:visual-multi}
\vspace{5pt} %
\scalebox{1}{
\begin{adjustbox}{max width=\textwidth}
\begin{tabular}{l p{12.5cm} }
\toprule
User & I want to have a pet, but I don't know whether to get a cat or a dog. Do you have any suggestions?\\
LLaDA & \includegraphics[width=\linewidth,valign=t]{imgs/visualization/multi0.jpg}\\
User & I'm very busy at work, so it seems like a cat would be a more suitable option. Do you have any suggestions for getting a cat? \\
LLaDA & \includegraphics[width=\linewidth,valign=t]{imgs/visualization/multi1.jpg}\\
\bottomrule
\end{tabular}
\end{adjustbox}
}
\end{minipage}
\end{table}



\begin{table}[t!]
  \begin{minipage}{0.99\textwidth}
\centering
\caption{\bf Single-turn Dialogue Cases of LLaDA 8B Instruct.}  
\label{tab:case-single}
\vspace{5pt} %
\scalebox{1}{
\begin{adjustbox}{max width=\textwidth}
\begin{tabular}{l p{12.5cm} }
\toprule
\emph{Daily life} & \\
~~~~User & What's a piece of advice you've received that you've found helpful?\\
~~~~LLaDA & One helpful piece of advice I've received is to always be open to learning and growth. I've found that being willing to explore new ideas and experiences has helped me become more resilient and adaptable. By embracing a growth mindset, I've been able to identify opportunities for improvement and development, both personally and professionally. \\

\midrule
\emph{Daily life} & \\
~~~~User & Please recommend me three famous movies\\
~~~~LLaDA &  Certainly! Here are three famous movies that have achieved widespread popularity and critical acclaim: \\
& 1. **The Shawshank Redemption** (1994) - Inspired by a novella and directed by Frank Darabont, this film tells the story of Andy Dufresne, a man convicted of murdering his wife, dealing with his harsh life in Shawshank Prison and forming an unlikely friendship with a fellow inmate. \\
& 2. **The Godfather** (1972) - Directed by Francis Ford Coppola, this film is based on Mario Puzo's novel of the same name. It follows the story of the Corleone family, particularly Michael Corleone, who rises to power and takes over the family's criminal empire. \\
& 3. **The Empire Strikes Back** (1980) - Directed by George Lucas, this film continues the story of a young Skywalker in the Star Wars saga. It follows Luke Skywalker as he helps the Rebel Alliance fight against the Galactic Empire and destroy the powerful weapon, the Death Star. \\
& All of these films are celebrated for their cinematic achievements, powerful storytelling, and memorable characters. \\

\midrule
\emph{Translation} & \\ 
~~~~User & Please help me translate into Chinese: `What is now proved was once only imagined' \\
~~~~LLaDA & \begin{CJK}{UTF8}{gbsn}现在被证明的曾经只是想象\end{CJK} \\

\midrule
\emph{Translation} & \\ 
~~~~User & Please help me translate into German: `What is now proved was once only imagined' \\
~~~~LLaDA & Wer heute bestätigt ist, war einst nur immaginiert. \\

\midrule
\emph{Math} & \\ 
~~~~User & Emma decides to run 4 sprints 5 times a week. She runs 75 meters each sprint. How many total meters does she run a week? \\
~~~~LLaDA & Emma runs 4 sprints each time, so she runs a total of 4 x 5 = 20 sprints a week. Each sprint is 75 meters long, so Emma runs a total of 20 x 75 = 1500 meters a week. The final result is 1500 \\

\midrule
\emph{Code} & \\ 
~~~~User & I am writing python code, I have an array, please help me sort it from largest to smallest \\
~~~~LLaDA &Certainly! You can sort an array from the largest to smallest in Python using the `sort()' method or the `sorted()` function. Here's an example:\\
    & \#\#\# Using `sort()' method:\\
    & my\_array = [3, 1, 4, 1, 5, 2] \\
    & my\_array.sort(reverse=True) \\
    & print(my\_array) \\
\bottomrule
\end{tabular}
\end{adjustbox}
}
\end{minipage}
\end{table}



\begin{table}[t!]
  \begin{minipage}{0.99\textwidth}
\centering
\caption{\bf Multi-turn Dialogue Cases of LLaDA 8B Instruct.}
\label{tab:case-multi}
\vspace{5pt} %
\scalebox{1}{
\begin{adjustbox}{max width=\textwidth}
\begin{tabular}{l p{12.5cm} }
\toprule
\emph{Daily life} & \\
~~~~User & How do I prepare for a job interview?\\
~~~~LLaDA & To prepare for a job interview, you should research the company and the job, practice common interview questions, and dress appropriately. \\
~~~~User & What are some common interview questions?\\
~~~~LLaDA & Some common interview questions include ``Tell me about yourself'' and ``Why do you want to work for us''?\\

\midrule
\emph{Daily life} & \\
~~~~User & How do I start learning programming?\\
~~~~LLaDA & To learn programming, you can start by choosing a programming language that you are interested in. Some popular languages include Python, Java, and C++. \\
~~~~User & Which programming language should I learn first?\\
~~~~LLaDA & It depends on your interests and goals. Python is a great language for beginners for its simplicity and versatility.\\
~~~~User & How long does it take to learn Python?\\
~~~~LLaDA & It can take anywhere from a few weeks to several months to learn Python.\\
~~~~User & What are some good Python projects for beginners?\\
~~~~LLaDA & Some good Python projects for beginners include building a simple calculator, a game, and a web application.\\

\midrule
\emph{Math} \\
~~~~User & Natalia sold clips to 48 of her friends in April, and then she sold half as many clips in May. How many clips did Natalia sell altogether in April and May? \\
~~~~LLaDA & Natalia sold clips to 48 friends in April. She sold half as many clips in May, so she sold 48/2 = 24 clips in May. Altogether, Natalia sold 48 + 24 = 72 clips in April and May. The final result is 72. \\ 
~~~~User & Natalia sold three times as many items in June as she did in May. How many clips did Natalia sell altogether in these three months? \\
~~~~LLaDA & Natalia sold clips to 48 friends in April. She sold half as many clips in May, so she sold 48/2 = 24 clips in May. Natalia sold three times as many items in June as she did in May, so she sold 24 * 3 = 72 clips in June. Altogether, Natalia sold 48 + 24 + 72 = 144 clips in these three months. The final result is 144. \\
\bottomrule
\end{tabular}
\end{adjustbox}
}

\end{minipage}
\end{table}



\begin{table}[t!]
  \begin{minipage}{0.99\textwidth}
\centering
\caption{\bf Poem Reversal Completion Cases where LLaDA 8B Instruct Succeeds but GPT-4o Fails.}  
\label{tab:case-reversal}
\vspace{5pt} %
\scalebox{1}{
\begin{adjustbox}{max width=\textwidth}
\begin{tabular}{l p{12.5cm} }
\toprule
\emph{Question1} & \\
~~~~User & \begin{CJK}{UTF8}{gbsn}`但闻人语响'的上一句是什么?\end{CJK}\\
~~~~GPT-4o & \begin{CJK}{UTF8}{gbsn}``但闻人语响''的上一句是``无可奈何花落去'',这两句出自唐代词人李清照的（武陵春，春晚: 无可奈何花落去，但闻人语响。这首词表达了作者在春天时节的孤独与感伤，借花落和人语声的描写，反映出词人内心的无亲与失落。
\end{CJK} \\
~~~~LLaDA & \begin{CJK}{UTF8}{gbsn}``但闻人语响''的上一句是``空山不见人''\end{CJK} \\
\midrule
\emph{Question2} & \\
~~~~User & \begin{CJK}{UTF8}{gbsn}`夜静春山空'的上一句是什么?\end{CJK}\\
~~~~GPT-4o & \begin{CJK}{UTF8}{gbsn}``夜静春山空''的上一句是``空山新雨后''.这两句出自唐代诗人王维的（山居秋膜）:空山新雨后，夜静春山空。这首诗描绘了春天山中雨后的寂静与空灵，表现出自然的宁静与和谐美。
\end{CJK} \\
~~~~LLaDA & \begin{CJK}{UTF8}{gbsn}``夜静春山空''的上一句是``人闲桂花落''\end{CJK} \\

\bottomrule
\end{tabular}
\end{adjustbox}
}
\end{minipage}
\end{table}


\begin{table}[t!]
    \centering 
    \caption{\textbf{Detailed results of LLaDA in Fig.~\ref{fig:scaling}.} "-" indicates missing values, which do not affect the observations regarding the scalability of LLaDA. These missing values are due to hardware failures. } 
    \label{tab:scaling_llada}
    \vspace{0.2cm}
    \begin{adjustbox}{max width=\textwidth}
    \begin{tabular}{ccc|cccccc}
    \toprule
    Model & Training Tokens & FLOPs & MMLU & CMMLU & ARC-C & PIQA & GSM8K & HumanEval \\
    \midrule
    LLaDA 1B & 37.75B & 2.20e20 & 25.52 & 25.95 & 25.17 & 59.41 & 1.82 & 0.00 \\
    LLaDA 1B & 88.08B & 5.13e20 & 27.11 & 26.52 & 26.96 & 61.86 & 3.03 & 1.83 \\
    LLaDA 1B & 138.41B & 8.06e20 & 29.32 & 27.13 & 30.20 & 63.38 & 2.35 & 0.00 \\
    LLaDA 1B & 239.08B & 1.39e21 & 31.48 & 30.77 & 27.99 & 63.11 & 3.26 & 1.22 \\
    LLaDA 1B & 352.32B & 2.05e21 & 35.86 & 34.35 & 31.31 & 65.34 & 3.64 & 3.05 \\
    LLaDA 1B & 461.37B & 2.69e21 & 31.86 & 30.98 & 30.12 & 65.51 & 2.35 & 0.61 \\
    LLaDA 8B & 62.91B & 2.63e21 & 32.22 & 28.5 & 30.20 & 63.82 & 3.87 & 2.44 \\
    LLaDA 8B & 125.83B & 5.27e21 & 33.39 & 33.9 & 34.64 & 66.54 & 8.72 & 3.66 \\
    LLaDA 8B & 251.66B & 1.05e22 & 42.84 & 40.59 & 40.10 & 69.04 & 15.31 & 3.66 \\
    LLaDA 8B & 377.49B & 1.58e22 & 45.11 & 43.99 & 39.25 & 68.61 & 25.40 & 9.76 \\
    LLaDA 8B & 503.32B & 2.11e22 & 43.57 & 41.38 & 42.06 & 70.24 & 27.52 & 9.76 \\
    LLaDA 8B & 629.14B & 2.63e22 & 48.80 & 47.13 & 42.24 & 72.09 & 30.10 & 12.80 \\
    LLaDA 8B & 679.48B & 2.85e22 & 49.61 & 48.19 & 41.30 & 70.84 & 26.31 & 8.54 \\
    LLaDA 8B & 792.72B & 3.31e22 & 50.88 & 49.01 & 42.58 & 70.51 & 31.99 & 6.10 \\
    LLaDA 8B & 981.47B & 4.11e22 & 49.47 & 48.10 & 40.27 & 71.38 & - & 6.10 \\
    LLaDA 8B & 1107.30B & 4.64e22 & 51.13 & 47.57 & 41.13 & 69.26 & 36.69 & 10.37 \\
    LLaDA 8B & 1233.13B & 5.16e22 & 50.52 & 49.72 & 45.05 & 71.49 & 38.97 & 9.76 \\
    LLaDA 8B & 1358.95B & 5.69e22 & 54.61 & 53.97 & 49.40 & 74.05 & 48.14 & 17.68 \\
    LLaDA 8B & 1547.70B & 6.48e22 & 57.38 & 56.04 & 49.49 & 74.59 & 53.30 & 20.73 \\
    LLaDA 8B & 1975.52B & 8.27e22 & 58.52 & 57.87 & 50.68 & 75.35 & - & 19.51 \\
    \bottomrule     
    \end{tabular}
    \end{adjustbox}
    \vspace{-.2cm}
\end{table}

\begin{table}[t!]
    \centering 
    \caption{\textbf{Detailed results of the autoregressive baelines in Fig.~\ref{fig:scaling}.}}
    \label{tab:scaling_arm}
    \vspace{0.2cm}
    \begin{adjustbox}{max width=\textwidth}
    \begin{tabular}{ccc|cccccc}
    \toprule
    Model & Training Tokens & FLOPs & MMLU & CMMLU & ARC-C & PIQA & GSM8K & HumanEval \\
    \midrule
    ARM 1B & 37.75B & 2.20e20 & 25.47 & 25.38 & 30.20 & 67.36 & 2.20 & 4.88 \\
    ARM 1B & 88.08B & 5.13e20 & 24.67 & 25.23 & 33.96 & 70.02 & 7.51 & 10.37 \\
    ARM 1B & 138.41B & 8.06e20 & 29.25 & 27.48 & 33.45 & 70.29 & 8.34 & 9.76 \\
    ARM 7B & 17.30B & 6.02e20 & 26.92 & 25.18 & 21.02 & 57.18 & 1.29 & 1.22 \\
    ARM 7B & 34.60B & 1.20e21 & 25.83 & 25.38 & 24.07 & 62.84 & 1.59 & 2.44 \\
    ARM 7B & 86.50B & 3.01e21 & 24.41 & 24.90 & 25.42 & 71.11 & 2.88 & 7.93 \\
    ARM 7B & 173.02B & 6.02e21 & 26.20 & 24.78 & 26.10 & 74.27 & 6.67 & 9.15 \\
    ARM 7B & 207.62B & 7.23e21 & 30.36 & 28.86 & 31.86 & 74.48 & 8.57 & 12.80 \\
    ARM 7B & 224.92B & 7.83e21 & 29.49 & 32.26 & 31.19 & 74.37 & 8.95 & 8.54 \\
    ARM 7B & 242.22B & 8.43e21 & 33.62 & 31.38 & 34.92 & 75.41 & 10.84 & 9.15 \\
    ARM 7B & 259.52B & 9.03e21 & 34.11 & 34.20 & 32.88 & 75.19 & 9.33 & 10.98 \\
    ARM 7B & 311.43B & 1.08e22 & 35.66 & 35.49 & 36.61 & 75.14 & 11.30 & 10.37 \\
    ARM 7B & 363.33B & 1.26e22 & 34.54 & 37.67 & 34.58 & 76.55 & 12.28 & 14.02 \\
    ARM 7B & 415.24B & 1.45e22 & 35.37 & 38.37 & 35.25 & 76.39 & 14.40 & 12.80 \\
    ARM 7B & 449.84B & 1.57e22 & 39.51 & 39.24 & 34.92 & 76.82 & 14.94 & 14.63 \\
    ARM 7B & 519.09B & 1.81e22 & 40.30 & 40.69 & 37.29 & 77.15 & 14.03 & 14.63 \\
    ARM 7B & 778.57B & 2.71e22 & 43.33 & 43.50 & 38.31 & 77.53 & 17.59 & 14.63 \\
    ARM 7B & 1038.09B & 3.61e22 & 45.06 & 46.12 & 41.69 & 77.86 & 20.02 & 15.85 \\
    ARM 7B & 1384.12B & 4.82e22 & 47.63 & 48.18 & 47.80 & 76.93 & 22.82 & 15.24 \\
    ARM 7B & 2076.18B & 7.23e22 & 47.68 & 50.85 & 44.07 & 77.37 & 24.79 & 14.63 \\
    ARM 7B & 2214.59B & 7.71e22 & 49.26 & 52.08 & 53.56 & 77.69 & 27.37 & 17.07 \\
    \bottomrule     
    \end{tabular}
    \end{adjustbox}
    \vspace{-.2cm}
\end{table}






\end{document}